\documentclass[11pt]{article}
\usepackage{amsmath,amssymb,amsthm}
\usepackage{fullpage}
\usepackage[hidelinks]{hyperref}
\newtheorem{theorem}{Theorem}
\newtheorem{hypothesis}{Hypothesis}
\newtheorem{conjecture}{Conjecture}
\newtheorem{corollary}{Corollary}
\newtheorem{lemma}{Lemma}
\newtheorem{question}{Question}
\theoremstyle{definition}
\newtheorem{definition}{Definition}
\theoremstyle{remark}
\newtheorem{remark}{Remark}
\newtheorem{namedthm}{Theorem}
\title{Causal Erasure: the Process-Rate Conditional Coordinate\\[2pt]
\large GO-14 statement v0.6 --- Theorems R, C and T (the convex
program, and the $n$-transfer lemma), the (H$^\ast$) \emph{repair}
--- the cross-block read $X$, the lemma chain L1--L4, and the
split-causal sub-family --- and T$^\ast$ v2, under the U-independence
hypothesis and the cross-read hypothesis (H$^{\ast\ast}$)}
\author{Geometric Observation program}
\date{v0.6 --- August 2026 (v0.5 superseded: the H-repair lemma chain
L1--L4 proves $D_1+D_2\le c(\Delta;m)+(2-\mathbf 1\{\Delta=0\})X$ for
\emph{every} record of $\mathcal F_0$, $X$ the cross-block read, using
no optimizer structure at all; consequently (H$^\ast$) is a
\emph{theorem} on the split-causal sub-family
$\mathcal F_0^{\mathrm{sc}}(m)$ with this document's own constants, and
is demoted to a \emph{corollary} of the strictly weaker,
better-margined, monotone-verified hypothesis (H$^{\ast\ast}$) on the
single scalar $X$; added: the $\Pi$-retraction lemma, the certified
$\phi^{\mathrm{sc}}$ brackets, the extended monotone $m\times\Delta$
table that retires the ``$m\in\{8,16\}$ only'' caveat, the repaired
per-block feasibility step of T$^\ast$(ii), the recomputed plateau
corollary $0.5514005>0.5479448$, and the \emph{named} obstruction that
replaces the old repair-path remark. v0.5 history: Theorem~T --- exact
subadditivity, Fekete existence of the limit and the
\emph{unconditional} lower transfer side --- with (H$^\ast$) carrying
the boundary-charged reverse direction and the executed within-class
certification of the diagonal ladder; six R-IND-5 transfer
restatements in force, including the refutation of ``$\le\kappa$ per
side'' as an $\mathcal F_0$ lemma. v0.4 history: Theorems~R and~C with
the certified anchor table and the sandwich-margin interval, seven
R-IND-5 restatements. v0.2 history: process-limit conjecture, hard
hypothesis, dated erratum to Theorem~1(v); the v0.1 sign error was
corrected at v0.2)}
\begin{document}
\maketitle

\begin{abstract}
GO-12 settled the single-record causal-prefix eraser (the Kalman
family); the process-rate object --- block records erased cell by
cell, each against only the causal prefix of an aging reference ---
was promoted to a first theorem at v0.2 (netted, GO-P-2026-076).
v0.3 added the process-limit face: the $n\to\infty$ limit of
the coordinate exists at $O(1/n)$, closes at the smoothing-pole rate
$\lambda_s^{2\Delta}$ with a constant that is \emph{not} matched by
any tested causal-conditioning spectrum, and is driven by an
explicit, testable record mechanism (horizon-matched
$V$-cancellation). This revision adds the lower-bound theory: in
moment coordinates $(H,\Gamma)$ the family optimization is a jointly
convex program --- Theorem~R (a schedule-general moment-coordinate
representation of $nL_a$) and Theorem~C (joint convexity, so every
KKT point is the global family minimum and strong duality yields
computable lower bounds). All seven full-space anchors are now
\emph{certified two-sided} (widths $3.45\times10^{-6}$--$1.7\times
10^{-5}$), and the sandwich margin is a certified interval:
$\min L_a(0)-\mathrm{block}_{16}\in[0.0317872,\,0.0317906]$.
This revision adds the \emph{$n$-transfer} theory. Theorem~T is
unconditional: block concatenation gives exact subadditivity
$f(n_1+n_2)\le f(n_1)+f(n_2)$, so by Fekete the process limit
$L^\infty=\lim_n\phi_n=\inf_n\phi_n$ \emph{exists}, and the lower
transfer side $0\le\phi_n-L^\infty$ is a theorem --- every certified
finite-$n$ value bounds the limit from above. The \emph{reverse}
direction is not: the boundary charge is carried by a numbered hard
hypothesis on dyadic-chain optimizers, because
``$\le\kappa$ per side'' is \emph{false} as a family statement
($D$-feasible $U$-independent counterexamples reach $26\kappa$; no
universal constant exists). The diagonal ladder is now certified
\emph{within class}.
This revision \emph{repairs} that hypothesis. Introducing the
\emph{cross-block read} $X=I(\hat Y^{b_1};W^{b_2}\mid W^{b_1})$, a
four-step lemma chain --- $\kappa$ is an \emph{identity} not a bound,
$I(E;T^{b_2})\le\kappa+X$, and the two per-side charge bounds --- gives
$D_1+D_2\le c(\Delta;m)+(2-\mathbf 1\{\Delta=0\})X$ for \emph{every}
record of $\mathcal F_0$, with no optimizer structure anywhere. Three
things follow. (a) The refutation is \emph{explained}: the
counterexamples are exactly the records with enormous $X$ ($18.46$ and
$29.27$ bits), and the bound $I(E;T^{b_2})\le\kappa+X$ is \emph{tight}
to $5.7\times10^{-7}$ at one of them. (b) (H$^\ast$) is a
\emph{theorem}, with this document's own constants, on the
split-causal sub-family $\mathcal F_0^{\mathrm{sc}}(m)$ --- a convex
linear section, so Theorem~C certifies it two-sided, at a
split-causality value gap that is flat in $n$ ($\le0.0044$ bits, and
$\approx0.0041$ at $n=16$), and onto which an explicit distortion- and
block-law-preserving retraction $\Pi$ maps every record. (c) The
transfer results
become conditional on (H$^{\ast\ast}$): \emph{one scalar}, $X\le\bar X$
at dyadic-chain optimizers, measured $6.3\times10^{-3}$ against an
admissible $0.117$ ($18\times$ headroom, against (H$^\ast$)'s
$5.5\times$) and monotone decreasing in $m$ over $m\in\{8,12,16,24\}$.
The $\Delta=0$ plateau survives the larger constant at
$0.5514005>0.5479448$ (margin $3.456\times10^{-3}$, $203\times$ the
quoted bracket width) and the base-$24$ negative control still fails,
now by $3.11\times10^{-4}$. A separate hole in T$^\ast$(ii), whose cost
would have exceeded the whole repair term, is closed by convexity of
the value function $D\mapsto\phi_m(D)$.
Three further unconditional results are added. \emph{Theorem~K}: the
KKT system of the convex program \emph{splits by column block},
because $E[S\mid W]=V$ and $Y$ are both coordinate blocks of $W$, so
the $Q$-factor and the distortion gradient have disjoint supports; it
gives $A_y=\theta N$, $N^{-1}=\theta I+V\Sigma_s^{-1}V^{\!\top}$,
$A_v=-NV\Sigma_s^{-1}V_S^{\!\top}$, hence the first $m$-uniform
regularity of the program ($0\prec N\preceq I/\theta$,
$0\prec A_y\preceq I$, $\theta\ge1.119$) and, as a \emph{corollary},
the transpose symmetry that v0.3 could only measure.
$\phi^{\mathrm{sc}}_{2m}\le\phi_m$ holds unconditionally, giving a
third and preferable hypothesis (H$^{\ast\ast\ast}$) on the
split-causality value gap --- preferable because that gap, unlike $X$,
is \emph{two-sided certifiable} by Theorem~C ($27\times$ headroom).
And the repair path itself is now executed rather than promised: an
$m$-uniform conversion theorem reduces (H$^{\ast\ast}$) to kernel-decay
envelopes plus a noise floor, yielding $\bar X=0.065038$ against an
admissible $0.116930$; what remains is named exactly --- the operator
whose inverse produces the optimizer is built \emph{from} the
optimizer, so inverse-closedness returns its own hypothesis, and a
fixed-point argument in a weighted Banach algebra is the missing step.
Two scoping results are load-bearing and
stated as hard terms: (i) the record family --- jointly Gaussian
with $(V,Y)$, \emph{independent of the reference noise $U$} --- is a
numbered hypothesis on the chain rule and on every minimum quoted
here (U-coupled records break the identity and invert the sandwich);
(ii) all numerical minima remain \emph{class-conditional or
finite-$n$ upper bounds}; nothing here relies on them being global.
\end{abstract}

\section{Setting}
Sources as in GO-12: $V$ AR(1) with pole $a$, unit stationary
variance; $Y_t=\rho V_t+N_t$ with $\operatorname{Var}(Y)=1$;
reference $S_t=V_t+U_t$, $\operatorname{Var}(U)=\tau^2$. A block
record $\hat Y^n$ of $Y^n$ (per-symbol distortion
$\tfrac1n\sum_tE(Y_t-\hat Y_t)^2\le D$) is erased cell by cell; the
eraser of cell $t$ holds $S^{t+\Delta}$ and the already-erased cells
$\hat Y^{t-1}$. The coordinate is
\[
L_a(\Delta)\;=\;\frac1n\sum_{t=1}^{n}
I\!\left(T^n;\hat Y_t\,\middle|\,S^{t+\Delta},\hat Y^{t-1}\right),
\qquad T=(V,Y).
\]

\begin{hypothesis}[U-independence --- hard, load-bearing]
\label{hyp:uind}
The record family $\mathcal F_0$: $\hat Y^n$ jointly Gaussian with
$(V^n,Y^n)$ and \emph{independent of $U^n$ given $(V^n,Y^n)$};
concretely $\hat Y=A_yY+A_vV+Z$ with $Z\perp(V,N,U)$ (no
$U$-coupling, $A_u=0$). Theorem~\ref{thm:main} (the chain rule and
every item), Conjecture~\ref{conj:proc}, and \emph{every minimum
quoted in this document} are stated under this hypothesis; no seal
of the process-limit face is valid without it.
\end{hypothesis}

\noindent\textbf{Extension and counter-values (R-IND-5, on record).}
The chain rule of Theorem~\ref{thm:main}(i) extends \emph{exactly}
to $U$-coupled records $\hat Y=A_yY+A_vV+A_uU+Z$ if and only if the
$U$-coupling is $\Delta$-lag causal, $A_u[t,s]=0$ for $s>t+\Delta$.
Outside that: the identity is \emph{false} (measured residuals
$0.015$--$0.40$ at random $U$-couplings; $4.59$ bits at an optimizer
endpoint), and the minima collapse --- the trivial $\Delta$-causal
$U$-record $\hat Y_t=0.54\,Y_t+0.23\,S_t+\mathcal N(0,0.1625)$ is
$D$-feasible and achieves $L_a(0)=0.539536<0.566758$ (the family
min), and a general $U$-coupled certified-feasible record achieves
$L_a(0)=0.092864<\mathrm{block}_{16}=0.534968$: the sandwich of
item~(iii) is \emph{inverted}. The $U$-coupled variant is a
different, far cheaper coordinate and is \emph{open}.

\section{The theorem}
Write $\lambda_s=a(1-K)$ for the steady-state Kalman
\emph{smoothing} pole ($K$ the steady Kalman gain of $S$ on $V$);
at the netted parameters $(a,\rho^2,\tau^2)=(0.8,0.49,0.4)$,
$\lambda_s=0.354554$, $\lambda_s^{-2}=7.955$.

\begin{theorem}[Causal erasure, first netting --- GO-P-2026-076]
\label{thm:main}
For every jointly Gaussian record on the window $t=1,\dots,n$
satisfying Hypothesis~\ref{hyp:uind}:
\begin{enumerate}
\item[(i)] \textbf{Exact chain rule.}
$n\,L_a(\Delta)=I(T^n;\hat Y^n\mid S^n)+C_\Delta$ with
\[
C_\Delta=\sum_t I\!\left(\hat Y_t;S_{>t+\Delta}\,\middle|\,
S^{t+\Delta},\hat Y^{t-1}\right)\;\ge\;0
\]
the \emph{smoothing-leakage charge}; $C_\Delta=0$ iff the record is
$\Delta$-lag causally simulatable from $S$. The identity is exact
for every record in $\mathcal F_0$ and every $\Delta$ (numerical
residual $\sim10^{-14}$); for the exact $U$-coupled extension and
its failure outside the family see the counter-values above.
\item[(ii)] \textbf{Collapse is a feasibility face, not an
impossibility.} Records that collapse ($C_\Delta=0$) while carrying
information exist --- e.g.\ any $V$-free ($N$-only) record --- so
universal non-collapse is \emph{false as first conjectured}. The
correct statement: for $\tau^2>0$, $\Delta\le n-2$ and
\[
D\;<\;\rho^2\,\frac{n-\Delta-1}{n},
\]
every $D$-feasible collapsing record must be $V$-free on all cells
$t\le n-\Delta-1$, which the distortion budget forbids; at
$D\ge\rho^2(n-\Delta-1)/n$ a feasible collapsing boundary record
exists (exhibited at $\Delta=6$, $n=16$, $D=0.3$). Strictness of
$\min L_a>\mathrm{block}/n$ in the sub-boundary regime routes
through block-optimality plus $C_\Delta>0$ at the (per-mode unique)
block optimizer --- not through universal non-collapse.
\item[(iii)] \textbf{Class-conditional sandwich, strictly
decreasing.} With $\mathrm{block}_n=\min\,\tfrac1nI(T^n;\hat
Y^n\mid S^n)$ and $\mathrm{slice}(\Delta)$ the netted single-letter
slice coordinate,
\[
\mathrm{block}_n\;<\;\min_{\text{records}} L_a(\Delta)\;\le\;
L_a^{\mathrm{diag}}(\Delta)\;<\;\mathrm{slice}(\Delta),
\]
where $L_a^{\mathrm{diag}}$ is the diagonal-class value; the upper
bounds are \emph{strictly decreasing in $\Delta$} (v0.1 said
increasing --- a sign error: more prefix access can only cheapen
causal erasure), and the lower margin closes geometrically at rate
$\lambda_s^{2\Delta}$ (overstatement of the diag-class bound over
the true minimum $\le1.07\times10^{-3}$ at $\Delta=2$,
$\le2.4\times10^{-6}$ at $\Delta=5$; $n=16$, $D=0.3$).
\item[(iv)] \textbf{The smoothing pole --- cell-local scope.} For
\emph{time-local} (cell-diagonal) records, the per-interior-cell
leakage ratio $\mathrm{leak}_t(\Delta)/\mathrm{leak}_t(\Delta+1)\to
\lambda_s^{-2}$ in the window interior; the \emph{aggregate}
per-symbol ratio carries the finite-window cell-count prefactor
$(n-\Delta-1)/(n-\Delta-2)$. The exponent is \emph{record-dependent}:
code-mixing records measurably decay slower (the reverse-water-filling
record shows a stable $\approx0.255$/lag, empirically
$\approx\lambda_s^2/\rho^2$ --- characterization open). No constant
measured ratio is quoted --- the pole is an asymptotic with a named
prefactor and a named scope.
\item[(v)] \textbf{Not bookkeeping.} $L_a$ is not the fixed-lag
static quadratic: the deficit of the bookkeeping candidate is
$n$-pinned, $\mathrm{gap}_n=\mathrm{static}(q_{\mathrm{path}})-
\mathrm{block}_n=0.0213/0.0263/0.0280$ at $n=8/16/24$, converging to
the spectral gain $\approx0.0313$
($\mathrm{block}_\infty=0.529950$ by per-frequency equal-slope
allocation; corrected value --- see the dated erratum,
Remark~\ref{rem:erratum}). Record memory and cross-cell code
structure carry real value.
\end{enumerate}
\end{theorem}

\begin{remark}[Class scoping --- load-bearing]
The diagonal (stationary-symmetric) class is \emph{provably
suboptimal} for $L_a$: at $\Delta=0$, $n=16$, $D=0.3$ the
first-order optimality certificate fails at the diagonal optimum
(feasible projected slope $8.6\times10^{-3}$ against
$10^{-6}$ noise) and explicit non-diagonal records achieve
$\min L_a(0)\le0.567353<0.572255$ (an improvement of
$4.9\times10^{-3}$ bits, two independent parameterizations agreeing
to $8\times10^{-5}$); the full-space search behind
Conjecture~\ref{conj:proc} has since sharpened the beat to
$\min L_a(0)\le0.5667581$ at $n=16$ (Remark~\ref{rem:anchors}).
Diagonal optimality \emph{remains proven} for the block coordinate.
Every quoted minimum is therefore a class-conditional or finite-$n$
upper bound inside the bracket of (iii); presenting $0.572255$ as
\emph{the} minimum would be an error.
\end{remark}

\begin{remark}[Rejected definition]
Definition (b) (memoryless eraser) is monotone but has the wrong
$\Delta\to\infty$ invariant (floors $0.0136$ bits above the netted
block coordinate) and is chain-rule-inconsistent; (c)
(Kostina--Hassibi causally-conditioned DI) collapses to the
full-path endpoint in Gaussian settings. Definition (a) is the
thermodynamically operational coordinate pending a reset-protocol
construction (open).
\end{remark}

\section{The convex program: Theorems R and C}
\label{sec:convex}
Status: lower-bound prover + independent R-IND-5 verification (own
evaluator, own polish, own moment box, fresh seeds; representation
identity to $2.0\times10^{-12}$ over 279 cells including edges; the
composition chain of Theorem~\ref{thm:cvx} re-proved independently;
$11{,}200$ verifier Jensen midpoints and $1{,}566$ curvature probes
with zero violations, on top of the prover's $32{,}000$). Harness:
\texttt{experiments/\allowbreak go14\_convexity.py}. Prereg 079
pending --- nothing in this section is sealed yet.

\emph{Moment coordinates.} Write $W=(V^n;Y^n)$ (so $T^n=W$),
$\Sigma_W=\operatorname{Cov}(W)$, and for a record $\hat Y=AW+Z$ of
the family $\mathcal F_0$ set
\[
H=\operatorname{Cov}(\hat Y,W)=A\Sigma_W\ \ (n\times2n),\qquad
\Gamma=\operatorname{Cov}(\hat Y)\ \ (n\times n).
\]
The family in these coordinates is the convex cone
\[
\mathcal D=\{(H,\Gamma):\ \Gamma-H\Sigma_W^{-1}H^{\!\top}\succeq0\}
\]
(the slack is $N_{\mathrm{cov}}=\operatorname{Cov}(Z)$), and the
distortion is \emph{affine}:
$n\,\mathrm{dist}=\operatorname{tr}\Sigma_Y
-2\operatorname{tr}(H\,\Sigma_W^{-1}\Sigma_{WY})
+\operatorname{tr}\Gamma$. Write $P=\Sigma_W^{-1}$ and
$Q=\Sigma_W^{-1}\Sigma_{WS}\Sigma_S^{-1}\Sigma_{WS}^{\!\top}
\Sigma_W^{-1}$.

\begin{namedthm}[R --- moment-coordinate representation,
schedule-general]
\label{thm:repr}
Assume Hypothesis~\ref{hyp:uind}. On the window $t=1,\dots,n$ let the
eraser of cell $t$ hold $S^{se(t)}$ for an \emph{arbitrary
nondecreasing access schedule} $se(\cdot)$, and define the
schedule-general counts
\[
k(j)\;=\;\#\{t:\ se(t)<j\},\qquad 1\le j\le n
\]
(the number of record cells already erased when $S_j$ enters the
interleaved order). Then for every $(H,\Gamma)\in
\operatorname{int}\mathcal D$:
\begin{align*}
2\ln2\cdot n\,L_a
&=\Bigl[\ln\det\bigl(\Gamma-HQH^{\!\top}\bigr)
      -\ln\det\bigl(\Gamma-HPH^{\!\top}\bigr)\Bigr]\\
&\quad+\sum_{j=1}^{n}\Bigl[\ln\operatorname{Var}(S_j\mid S^{j-1})
      -\ln\operatorname{Var}\bigl(S_j\mid S^{j-1},\hat Y^{k(j)}
      \bigr)\Bigr].
\end{align*}
\emph{Staircase reduction (explicit).} For the coordinate's own
schedule $se(t)=\min(t+\Delta,n)$ the counts reduce to
$k(j)=\max(0,\,j-\Delta-1)$. The reduced form is
\emph{prefix-staircase-specific}: on non-staircase nondecreasing
schedules only the general $k(j)=\#\{t: se(t)<j\}$ is correct ---
blind use of the staircase formula deviates by up to $1.05$ bits
(verifier-measured), and off-by-one variants of $k$ are detected at
$0.1$--$0.7$ bits.
\end{namedthm}

\noindent\emph{Proof skeleton.} (1) Factor the joint
$\ln\det$ of $(S^n,\hat Y^n)$ along the interleaved order twice:
collecting the $\hat Y$-pivots gives $\sum_t\ln\operatorname{Var}
(\hat Y_t\mid S^{se(t)},\hat Y^{t-1})$; collecting the $S$-pivots
gives $\sum_j\ln\operatorname{Var}(S_j\mid S^{j-1},\hat Y^{k(j)})$
--- this regrouping is the whole trick (and the reason single-cell
convexity is never needed). (2) U-independence enters \emph{exactly
once} as a premise --- the residuals $(U,Z)$ are independent given
$W$ --- but is load-bearing in \emph{two} legs of the computation:
$\operatorname{Cov}(S,\hat Y)=KH^{\!\top}$ with
$K=\Sigma_{WS}^{\!\top}\Sigma_W^{-1}$ (numerator/leak leg, giving the
$Q$-Schur complement $\Gamma-HQH^{\!\top}$ since
$Q=K^{\!\top}\Sigma_S^{-1}K$), and
$\ln\det\operatorname{Cov}(\hat Y\mid W,S^{\bullet})
=\ln\det(\Gamma-HPH^{\!\top})$ (denominator leg: conditioning on $S$
adds nothing given $W$). (3) The unconditional $S$-pivots telescope
against $\ln\det\Sigma_S$. This re-proves the 076 chain rule of
Theorem~\ref{thm:main}(i) in one line: the bracket is the block
term $I(T^n;\hat Y^n\mid S^n)$ and the $j$-sum is the leakage charge
$C_\Delta$ in its $S$-side (reversed) form. Verification: residual
$1.15\times10^{-12}$ over 144 prover cells and $2.0\times10^{-12}$
over 279 verifier cells including the edges
$\Delta\in\{0,n-2,n-1,n\}$ and random non-staircase schedules.\hfill
$\square$

\begin{remark}[U-independence: one premise, two legs --- and the
non-conflation rule]
\label{rem:ulegs}
Correcting only one leg does \emph{not} restore the identity: at
fixed dense-$U$-coupled records ($n=12$) the leg-by-leg
counter-values are $0.32$--$2.01$ bits (as-stated, numerator-only-
corrected, and denominator-only-corrected all break; the harness's
pinned record measures $2.68/3.51/0.83$ bits respectively), while
correcting \emph{both} legs restores exactness for arbitrary $A_u$
including anticausal couplings (residual at machine precision).
Non-conflation: $\Delta$-lag-causal $U$-coupling
($A_u[t,s]=0$ for $s>t+\Delta$) extends the \emph{record-space}
chain rule of Theorem~\ref{thm:main}(i) exactly, but does
\emph{not} rescue the moment-form Theorem~R as stated --- measured
moment-form residuals $2.4$--$3.5$ bits at $\Delta$-causal
couplings (harness pinned record: chain-rule residual
$2\times10^{-16}$, moment-form residual $3.22$ bits at $\Delta=2$).
The two extension results must not be conflated.
\end{remark}

\begin{namedthm}[C --- joint convexity; KKT points are global;
strong duality]
\label{thm:cvx}
On $\mathcal D$: (i) the block bracket
$B(H,\Gamma)=\ln\det(\Gamma-HQH^{\!\top})
-\ln\det(\Gamma-HPH^{\!\top})$ equals $-\ln\det(I-Z)$ with
$Z=R^{1/2}H^{\!\top}(\Gamma-HQH^{\!\top})^{-1}HR^{1/2}$ and
\[
R\;=\;P-Q\;=\;\Sigma_W^{-1}\Sigma_{W\mid S}\Sigma_W^{-1}
\;\succeq\;0 ;
\]
under the mapping $C\leftrightarrow H^{\!\top}$,
$V\leftrightarrow\Gamma$ this is \emph{exactly} the 074 convexity
lemma's $G$-form (GO-13 Theorem~5, block Schur lift), so $Z$ is
jointly matrix-convex on $\mathcal D$ with $Z\prec I$ on
$\operatorname{int}\mathcal D$ (Sylvester:
$\det(I-Z)=\det M_P/\det M_Q$), and $B$ is jointly convex. (ii) Each
leak term $-\ln s_j$ is convex: the pivot
$s_j=\operatorname{Var}(S_j\mid S^{j-1},\hat Y^{k(j)})$ is a Schur
pivot of a joint covariance that is \emph{affine} in $(H,\Gamma)$,
hence an infimum of affine functions of $(H,\Gamma)$ ---
concave --- and bounded below by $\tau^2>0$ (the eraser never sees
$U_j$ through a $U$-independent record). (iii) The distortion is
affine and $\mathcal D$ is a convex cone. Hence $n\,L_a$ is jointly
convex on $\mathcal D$; every KKT point of the
distortion-constrained program is the \emph{global} family minimum;
Slater's condition holds (interior scalar records), so strong
duality makes lower bounds computable: any $\mu\ge0$ and any point
$x_p$ give the certificate $p^\ast\ge\phi_\mu(x_p)-
\|\nabla\phi_\mu(x_p)\|\,R_{\mathrm{box}}$, with
$R_{\mathrm{box}}$ from the per-cell moment bounds
$\Gamma_{tt}\le(1+\sqrt{nD})^2$, $|H_{tj}|\le\sqrt{\Gamma_{tt}}$
(verifier: the box is valid and nearly tight --- an adversarial
one-cell record reaches $10.1748$ against the box's $10.1818$).
\end{namedthm}

\noindent\emph{Proof skeleton.} (i) is the 074 lift verbatim after
the transpose-convention check ($C^{\!\top}\!PC$ with $C=H^{\!\top}$
is $HPH^{\!\top}$ --- mapping exact to $0.0$); $R\succeq0$ by the
displayed closed form (verified to $1.7\times10^{-14}$;
$\operatorname{eigmin}(Q)\ge-10^{-15}$, PSD to machine). (ii) is
inf-of-affine concavity plus $-\ln$ composition on
$[\tau^2,\infty)$. (iii) is inspection. Numerics on record: zero
Jensen-midpoint violations (prover $32{,}000$ + verifier
$11{,}200$ + $1{,}566$ curvature probes), tangent plane at the
$(16,0)$ winner $400/400$ with min slack $+1.13\times10^{-2}$,
analytic gradient vs FD $1.2\times10^{-10}$.\hfill$\square$

\begin{remark}[Per-cell convexity is open --- the regrouping is the
proof]
\label{rem:percell}
Theorem~\ref{thm:cvx} proves convexity of the \emph{regrouped}
pieces (block bracket, $S$-side leak sum). Whether the individual
per-cell CMI terms of $L_a$ are themselves convex in $(H,\Gamma)$
remains \emph{open}: no violation was found in any sampling (prover
and verifier both, zero hits), but no proof exists either. Nothing
in this section relies on per-cell convexity.
\end{remark}

\begin{remark}[The diagonal class is a convex linear section ---
its ladder is UB-only because class $<$ family]
\label{rem:diagsection}
In $(H,\Gamma)$ the diagonal (stationary-symmetric) class is a
\emph{convex (linear) section} of $\mathcal D$: midpoints of
diagonal records stay in class (verifier: off-diagonal mass in the
$C_v$ eigenbasis at $(H,\Gamma)$-midpoints is $0$ to machine
precision). The earlier ``non-convex slice'' wording is refuted as
worded --- what is non-convex is the $(g,h,z)$ \emph{parameter
chart} (explicit midpoint witness on record), a different
statement. Consequently the diagonal ladder of
\S\ref{sec:proc} is upper-bound-only \emph{because the class is a
strict subset of the family} (Remark on class scoping,
\S2), not because of any slice non-convexity; and within-class
values are certifiable by the same Lagrangian machinery restricted
to the section. \textbf{That upgrade path has since been executed}:
all nine ladder cells and three block values are now certified
two-sided \emph{class} values (\S\ref{sec:ladder}). Certifying the
class does not certify the family --- the ladder remains an upper
bound on $\min L_a$ because class $<$ family.
\end{remark}

\medskip\noindent\textbf{Certificates on record (per-$n$).}
$\operatorname{eigmin}(P-Q)=1.0382\times10^{-2}\,/\,
5.5944\times10^{-3}\,/\,4.4309\times10^{-3}\,/\,
3.9529\times10^{-3}$ at $n=8/16/24/32$ --- \emph{decreasing in
$n$}; quoted as ``$\ge3.9\times10^{-3}$ for $n\le32$'' (no
$n$-uniform floor is claimed; the earlier single-number quote is an
erratum, corrected here). The noise-floor parameterization is
inactive at every winner and every polished point.

\subsection{Theorem K: the KKT system splits by column block}
\label{sec:K}
Two \emph{exact} structural facts of this model drive everything in
this subsection, and both are checkable in one line:
\[
K=\Sigma_{WS}^{\!\top}\Sigma_W^{-1}=[\,I\ \ 0\,],
\qquad
\Sigma_W^{-1}\Sigma_{WY}=[\,0;\,I\,]
\]
(verified to $2\times10^{-15}$): the first because $E[S\mid W]=V$ and
$V$ \emph{is} a coordinate block of $W$, so every $Q$-factor term lives
only on the $V$-columns of $A=H\Sigma_W^{-1}=[A_v,A_y]$; the second
because $Y$ is a coordinate block of $W$, so the distortion gradient
lives only on the $Y$-columns. \textbf{The two supports are disjoint},
so the stationarity condition splits by column block.

\begin{namedthm}[K --- KKT structure; \emph{unconditional},
$\mathcal F_0$, any $\Delta$, any nondecreasing schedule]
\label{thm:K}
Let $(H,\Gamma)$ be a KKT point of the distortion-constrained program
with multiplier $\mu$, and write $N=\Gamma-H\Sigma_W^{-1}H^{\!\top}
=\operatorname{Cov}(Z)$, $\theta=2\mu\ln2$, and
$(v_t,v^S_t,\sigma_t)$ for the \emph{record} pivots of the interleaved
order (cell $t$ predicted from $S^{se(t)},\hat Y^{t-1}$), collected
into $V\Sigma_s^{-1}V^{\!\top}$ and $V\Sigma_s^{-1}V_S^{\!\top}$. Then
\[
\text{(K1)}\ \ A_y=\theta N,\qquad
\text{(K2)}\ \ N^{-1}=\theta I+V\Sigma_s^{-1}V^{\!\top},\qquad
\text{(K3)}\ \ A_v=-N\,V\Sigma_s^{-1}V_S^{\!\top}.
\]
Verified at every certified optimizer --- relative residuals
$\le3.2\times10^{-8}$ at $(16,0)$, $(24,0)$, $(32,0)$, $(48,0)$,
$(16,1)$, $(16,2)$, $(24,1)$ (prover), and $\le8.3\times10^{-8}$ over
the six of those the harness recomputes from cold starts. The
residuals track the polish accuracy of the optimizer, not the
identities, which are exact; the harness gates them at $10^{-6}$ for
that reason.
\end{namedthm}

\noindent\emph{Corollaries, unconditional and $m$-uniform.}
\begin{enumerate}
\item[(a)] $A_y=\theta N$ is \emph{symmetric positive definite}.
\textbf{This upgrades the 078-era ``mechanism'' finding --- $A_y$ and
the noise kernel transpose-symmetric, all causal asymmetry confined to
$A_v$ --- from an empirical observation to a corollary of
Theorem~\ref{thm:K}} (measured symmetry defect
$1.2\times10^{-9}$--$1.4\times10^{-8}$, against the
$5\times10^{-5}$ at which it was first reported; see
Conjecture~\ref{conj:proc}(4)). And (K3) carries the causal asymmetry
through $V_S$, whose column support ends \emph{exactly} at $t+\Delta$
--- the structural candidate for the observed sign change at lag
$\Delta$. \emph{This is a lead, not a claim}: no proof that the sign
change must occur there is offered here.
\item[(b)] $0\prec N\preceq I/\theta$ and $0\prec A_y\preceq I$,
\emph{uniformly in $m$} --- the first $m$-uniform optimizer-regularity
statement in the program.
\item[(c)] $\theta\ge2\ln2\,\phi_n(D)/(1-D)\ge1.119$ (convexity of
$D\mapsto\phi_n(D)$, Remark~\ref{rem:feas}, together with
$\phi_n(1)=0$), hence $N\preceq0.894\,I$ unconditionally.
\end{enumerate}

\noindent\emph{Proof skeleton.} Stationarity in $\Gamma$ reads
$M_Q^{-1}-N^{-1}+\partial_\Gamma\bigl(-\sum_j\ln s_j\bigr)=-\theta I$.
The Cholesky decomposition of $J^{-1}$ along the interleaved order,
$J=\operatorname{Cov}(S^n,\hat Y^n)$, splits the pivot Gram into a
reference half and a record half summing to $J^{-1}$ (verified to
$6\times10^{-16}$); since $(J^{-1})_{\hat Y\hat Y}=M_Q^{-1}$, the
$M_Q^{-1}$ and reference-pivot terms combine into the record-pivot Gram
--- which is (K2). Stationarity in $H$ then splits by the disjoint
supports above: the $Y$-columns give $N^{-1}A_y=\theta I$, i.e.\ (K1);
the $V$-columns give (K3). \emph{Caution on reading (K2) and (K3)}:
the Gram must be the \emph{record}-pivot one --- substituting the
reference-pivot Gram makes (K2) false at the $47\%$ level and (K3)
false at the $120\%$ level.\hfill$\square$

\subsection{Certified two-sided anchors}
\label{sec:anchors}
Lagrangian polish + moment-box tangent bound
(Theorem~\ref{thm:cvx}), prover (P) and independent verifier (V)
each end-to-end; per anchor the tighter certified width of the two
is cited. Upper ends are the (feasibility-projected) winner values.

\begin{center}
\begin{tabular}{llll}
$(n,\Delta)$ & certified bracket for $p^\ast$ & width & source\\
\hline
$(16,0)$ & $[0.566754682,\ 0.5667581350]$ & $3.45\times10^{-6}$ & V\\
$(16,1)$ & $[0.5408029,\ \ \,0.5408064658]$ & $3.6\times10^{-6}$ & P\\
$(16,2)$ & $[0.5358901,\ \ \,0.5358939352]$ & $3.8\times10^{-6}$ & P\\
$(24,0)$ & $[0.5654096,\ \ \,0.5654138570]$ & $4.3\times10^{-6}$ & P\\
$(24,1)$ & $[0.5393228,\ \ \,0.5393377781]$ & $1.5\times10^{-5}$ & P\\
$(24,2)$ & $[0.5342628,\ \ \,0.5342788038]$ & $1.6\times10^{-5}$ & P\\
$(32,0)$ & $[0.5647249,\ \ \,0.5647418676]$ & $1.7\times10^{-5}$ & P\\
\end{tabular}
\end{center}

\noindent The 078 anchors are thereby \emph{certified family
values}, no longer upper bounds alone; minimizer \emph{uniqueness}
is still not claimed (no flat was found --- three polishes agree to
$2.8\times10^{-14}$ in value and $2\times10^{-6}$ in coordinates ---
but gauge flats are not excluded; ``an optimizer'', never
``the''). The harness reproduces four of the brackets end-to-end
from cold starts by $\mu$-bisection (convexity makes warm starts a
convenience, not a dependence), at widths $4.0/3.0/3.8/9.2\times
10^{-6}$ for $(16,0)/(16,1)/(24,0)/(32,0)$ and a
$\mathrm{block}_{16}$ bracket of width $1.7\times10^{-6}$.

\begin{corollary}[Certified sandwich margin at $n=16$, $D=0.3$]
\label{cor:margin}
\[
\min L_a(0)-\mathrm{block}_{16}\;\in\;[0.0317872,\ 0.0317906]
\]
(verifier-tightened interval). The block side inherits the sealed
$\mathrm{block}_{16}$ precision through the \emph{proven} diagonal
optimality of the block coordinate.
\end{corollary}

\begin{remark}[Feasible projection of winner upper ends]
\label{rem:proj}
The stored winner endpoints are distortion-infeasible by
$\sim4\times10^{-11}$ (raw optimizer termination). The one-line
feasible projection (blend toward the near-copy record) corrects
the upper ends by $\sim10^{-10}$ --- immaterial at the
$3.5\times10^{-6}$ width scale; all upper ends above are to be read
through this remark.
\end{remark}

\begin{remark}[``Certified'' means floating-point certificates]
\label{rem:fpcert}
All brackets here are \emph{floating-point} certificates: the
weak-duality and tangent-bound inequalities are exact in real
arithmetic and are evaluated in IEEE double precision without
interval arithmetic. The certified widths exceed the $f64$
evaluation error by roughly seven orders of magnitude; no claim of
formally verified arithmetic is made.
\end{remark}

\section{The process limit}
\label{sec:proc}
Everything in this section is under Hypothesis~\ref{hyp:uind} and at
the netted parameters, $D=0.3$. Status: probe + independent R-IND-5
verification (two evaluator routes agreeing to $2\times10^{-10}$;
all sealed 076 anchors reproduced); pre-registration (078) pending
--- no numerical face here is sealed yet.

\begin{conjecture}[Process limit of the causal-erasure coordinate ---
v0.2, R-IND-5-scoped]
\label{conj:proc}
Let $L^{\mathrm{fs}}_n(\Delta)=\min L_a(\Delta)$ over the full
(non-diagonal) family $\mathcal F_0$ on the window of length $n$,
and $L^{\mathrm{diag}}_n(\Delta)$ the diagonal-class value.
\begin{enumerate}
\item[(1)] \textbf{Existence at $O(1/n)$.} Both coordinates converge
as $n\to\infty$ with $O(1/n)$ finite-window drift (model attack:
fitted exponent $\alpha=0.996$--$1.001$). Two-point $1/n$ limits,
calibrated on the block face ($\le5\times10^{-6}$): full space
$L^{\mathrm{fs}}_\infty(\Delta)=0.562725/0.536400/0.531049$ at
$\Delta=0/1/2$ (three-$n$ consistency at $\Delta=0$:
$0.5627253/0.5627259/0.5627265$ from $n=16/24/32$ pairs); diagonal
class $0.568571/0.537401/0.531206/0.529977$ at $\Delta=0/1/2/4$
($\pm\le1.2\times10^{-5}$).
\item[(2)] \textbf{Closure rate, constant unclaimed.} The excess
$E_\infty(\Delta)=L^{\mathrm{diag}}_\infty(\Delta)-
\mathrm{block}_\infty$ obeys $E_\infty(\Delta)\sim
c\,\lambda_s^{2\Delta}$ \emph{as a $\Delta\to\infty$ asymptotic
only}: the per-lag ratio rises toward $\lambda_s^{-2}=7.955$
\emph{from below} ($5.18\to5.93\to6.87\to7.65\to$ within error at
$\Delta=5\to6$) and is \emph{unresolved beyond $\Delta\approx6$};
the constant is quoted as the bracket
$c_{\mathrm{diag}}\in[0.105,0.125]$ \emph{only} ---
$c(\Delta)=0.107/0.112/0.117$ at $\Delta=4/5/6$, still rising,
\emph{not converged through $\Delta=6$}. The full-family constant is
data-supported only as $c_{\mathrm{fs}}\in(0.07,0.125]$. No
single-letter identification of $c$ is claimed.
\item[(3)] \textbf{Strictly above every tested causal-conditioning
spectrum --- certified at every computed $n$; in the limit
conditional on (H$^\ast$) at $\Delta=0$ and on the $O(1/n)$ model at
$\Delta=1,2$.} At every computed anchor
($n\in\{16,24,32\}$, $\Delta\in\{0,1,2\}$) the \emph{certified lower
bracket end} of \S\ref{sec:anchors} sits strictly above the
causal-spectral allocation value (Remark~\ref{rem:refcon}) --- a
certificate-strict gap (smallest exceedance $\approx4\times10^{-3}$
bits at $(24,2)$, $\ge250\times$ the bracket width there), not an
optimizer report. In the limit, the plateaus
$+0.014781/+0.004057/+0.000795$ ($\Delta=0/1/2$) above the
allocation are $160$--$3000\times$ the calibrated extrapolation
error and carry the wrong shape (drifting per-lag rate, not the
allocation's clean $\lambda$-law). At $\Delta=1,2$ the limit
statement remains \emph{conditional on the $O(1/n)$ finite-window
model}. At $\Delta=0$ it is \emph{no longer} model-conditional:
Corollary~\ref{cor:plateau} derives
$L^\infty(0)\ge0.5514005>0.5479448$ from the certified $n=32$
bracket and the repaired transfer constant, \emph{conditional instead
on Hypothesis~\ref{hyp:hstarstar} \textup{(H$^{\ast\ast}$)}} --- or on
(D)$+$(F), at the smaller margin recorded in
Corollary~\ref{cor:plateau}. The subadditivity
half of the old $n$-monotonicity face is now closed unconditionally
(Theorem~\ref{thm:transfer}); what remains open is the cross-block
read at optimizers (Question~\ref{q:open}). Both closed-form
candidates
($c_{\mathrm{spec}}=0.0190\pm0.0007$,
$c_{\mathrm{static}}=0.026668$) are too small at every $\Delta$. No
claim is made that no single-letter form exists. The gap is genuine
cross-cell code value: the winner cuts leakage $0.033\to0.024$ by
paying block penalty the per-frequency relaxations cannot express.
\item[(4)] \textbf{Mechanism --- \emph{the symmetry half is no longer
a conjecture}.} At every full-space optimizer computed
($n=16/24/32$, $\Delta=0/1/2$; seven adversarial
starts converge to the same record, KKT certificate
$8.4\times10^{-6}$): the $Y$-coupling and noise kernels are
time-symmetric, where \emph{time-symmetric means transpose symmetry
$A_y=A_y^{\!\top}$} (relative residual $<5\times10^{-5}$ as first
measured; the stronger persymmetry $JA_y^{\!\top}J=A_y$ \emph{fails}
at the $\approx6\%$ level at the window edges --- transpose symmetry is
the invariant); \emph{all} causal structure lives in the $V$-coupling
$A_v$, whose interior-row kernel \emph{changes sign exactly at lag
$\Delta$} --- horizon-matched $V$-cancellation: the record aligns
$V$-content with the span the eraser will hold and anti-correlates
$V$ beyond the access horizon, suppressing smoothing leakage. This
is exactly what the symmetric-by-class diagonal family cannot
express (and explains the 076 beat).
\textbf{Status upgrade (v0.6).} The first half of this item is
\emph{no longer conjectural}: $A_y=\theta N$ is symmetric positive
definite at every KKT point by Theorem~\ref{thm:K}(K1), so the
transpose symmetry of $A_y$ and of the noise kernel, and the
confinement of all causal asymmetry to $A_v$, are \emph{corollaries},
not measurements (measured symmetry defect now
$1.2\times10^{-9}$--$1.4\times10^{-8}$, i.e.\ at polish accuracy). What
remains conjectural is the \emph{sign change exactly at lag $\Delta$};
Theorem~\ref{thm:K}(K3) supplies a structural candidate --- $A_v$
inherits its asymmetry from $V_S$, whose column support ends exactly at
$t+\Delta$ --- but no proof. Falsifiable per-cell at any $n$.
\end{enumerate}
\end{conjecture}

\begin{remark}[Reproducible anchors --- raw finite-$n$ values only]
\label{rem:anchors}
The anchors of record are \emph{raw finite-$n$ $L_a$ values}: family
min $L_a(0)\le0.5667581$ at $n=16$ (beats the 076-recorded
$0.567353$ by $5.9\times10^{-4}$; all start basins converge to it
within $10^{-7}$); full-space $L_a(0)\le0.56474187$ at $n=32$; the
diagonal ladder value $L_a^{\mathrm{diag}}(n{=}48,\Delta{=}6)
=0.5316222730$. The full-space upper ends are now \emph{certified
two-sided} (\S\ref{sec:anchors}); the diagonal-ladder values are now
certified two-sided \emph{within the class} (\S\ref{sec:ladder}) and
remain class-conditional upper bounds on the family minimum
(Remark~\ref{rem:diagsection}). The extrapolated $E_\infty(6)$ is quoted \emph{only}
with an explicit error bar --- $E_\infty(6)\approx4.4\times10^{-7}
\pm10\%$ (two-point pair sensitivity); it is \emph{not} a four-digit
anchor, the raw $L_a(48,6)$ above is.
\end{remark}

\begin{remark}[The causal-spectral allocation is a reference
construction]
\label{rem:refcon}
The causal-spectral allocation: equal-slope per-frequency
conditional quadratics with conditioning spectrum the $\Delta$-lag
causal Wiener error spectrum $S_e^{(\Delta)}$. Values
$0.547945/0.532344/0.530253$ at $\Delta=0/1/2$, verified to seven
digits by an independent (time-domain) route; it reproduces
$\mathrm{block}_\infty$ at the two-sided endpoint to
$1.6\times10^{-10}$, and its own excess obeys the
$\lambda_s^{2\Delta}$ law cleanly. It is a \emph{reference
construction only}: it carries no achievability map (not an upper
bound on $\min L_a$) and no data-processing direction (not a lower
bound). Item (3) of Conjecture~\ref{conj:proc} is a statement about
mismatch with the tested spectra, never a bound.
\end{remark}

\begin{remark}[Erratum --- dated 2026-08-05]
\label{rem:erratum}
v0.2 of this document quoted $\mathrm{block}_\infty=0.52991$ in
Theorem~\ref{thm:main}(v), propagated from the pass-1 coarse
frequency waterfilling; the converged per-frequency equal-slope
value is $\mathrm{block}_\infty=0.529950$ (corrected in item (v)
above; the spectral gain $\approx0.0313$ is unaffected). Same date,
probe-record erratum: the process-limit probe's diagonal excess at
$\Delta=4$ is $2.66\times10^{-5}$, not $2.75\times10^{-5}$.

\emph{Dated 2026-08-06 (transfer note).} Three further corrections,
all carried into \S\ref{sec:transfer}: (a) the constant was recorded
as $c(1)\le1.1257$; the computed value is $1.1257294$, so the
recorded bound is false in the fifth decimal and $c(1)\le1.1258$ is
used throughout (Remark~\ref{rem:const}). (b) The transfer note cites
a value $0.564729705$ for the certified lower end at $(32,0)$; no
computation on record produces it. The sealed value is
$\mathrm{LB}(32,0)=0.5647327869$
(\texttt{results/GO14-convexity.json}, \texttt{s5\_32\_0}), and it is
the only one used in Corollary~\ref{cor:plateau}; $0.564729705$ must
not be cited. (c) The plateau margin is $215\times$ the quoted
bracket width, not the $\approx250\times$ of the transfer note.
\end{remark}

\section{The $n$-transfer lemma: Theorem T, and the repair of
(H$^\ast$)}
\label{sec:transfer}
Status: prover + independent R-IND-5 verification (own \texttt{slogdet}
evaluator, own certified anchors, own re-derivation of the constants
and of the corollary arithmetic; six mandatory restatements, all in
force in this revision). Harness
\texttt{experiments/\allowbreak go14\_transfer.py}. Registration 080
pending --- nothing in this section is sealed. The repair material of
\S\ref{sec:repair} (the lemma chain, the split-causal sub-family, the
$\Pi$-retraction, the monotone table, the feasibility repair and the
recomputed corollary) and Theorem~\ref{thm:K} of \S\ref{sec:K} are
newer still: they carry their own harness,
\texttt{experiments/\allowbreak go14\_hrepair.py}, and await
registration 081. \textbf{Nothing in this revision contradicts a gate
of 080}: every 080 gate is either untouched or --- in the case of the
$\kappa$-per-side refutation --- \emph{explained} rather than
overturned.

Write $\phi_n(\Delta)=\inf_{\mathcal F_0}L_a(\Delta)$ for the family
value on the window of length $n$ (certified two-sided in
\S\ref{sec:anchors}) and $f(n)=n\,\phi_n(\Delta)$. Let
\[
\kappa\;=\;\tfrac12\log_2\frac{1}{1-a^2}\;=\;0.736966\ \text{bits}
\]
be the AR(1) boundary information at the netted pole $a=0.8$. For a
split of the $n$-window into blocks $b_1,b_2$ define the \emph{per-side
boundary charges}
\[
D_i\;=\;\sum_{t\in b_i}\Bigl[
I\bigl(T^{b_i};\hat Y_t\mid \mathrm{own}C_t\bigr)
-I\bigl(T^{b_i};\hat Y_t\mid \mathrm{big}C_t\bigr)\Bigr]\;\ge\;0,
\]
where $\mathrm{own}C_t$ is the conditioning set of cell $t$ inside its
own block and $\mathrm{big}C_t=S^{se(t)},\hat Y^{t-1}$ is the
big-window one; $D_i$ is exactly what the split loses at the boundary.
Write $E=(S^{b_1},\hat Y^{b_1})$ for the block-1 record-and-reference
history that the block-2 cells see in $\mathrm{big}C$ but not in
$\mathrm{own}C$.

\begin{namedthm}[T --- exact subadditivity, Fekete existence, and the
lower transfer side; \emph{unconditional}]
\label{thm:transfer}
Fix $\Delta$ and assume Hypothesis~\ref{hyp:uind}.
\begin{enumerate}
\item[(i)] \textbf{Exact subadditivity.}
$f(n_1+n_2)\le f(n_1)+f(n_2)$ for all $n_1,n_2\ge1$. Take
$\varepsilon$-optimal records of the two windows, independent of each
other, and concatenate: the result is a feasible record of the
$(n_1+n_2)$-window (the distortion constraint is the cell average),
and \emph{at every cell} the big-window per-cell term is at most the
own-window term. The numerator
$\operatorname{Var}(\hat Y_t\mid S^{se(t)},\hat Y^{t-1})$ can only
shrink when the conditioning set grows from the own block to the whole
window; the denominator
$\operatorname{Var}(\hat Y_t\mid T^{n},S^{se(t)},\hat Y^{t-1})$ is
\emph{unchanged}, because a block-diagonal record makes cell $t$
conditionally independent of the other block given its own block's
$T$ --- verified as an exact equality (relative deviation
$2\times10^{-14}$ over all cells, including the unequal split
$8+4\to12$). Letting $\varepsilon\downarrow0$ gives the inequality
for the infima.
\item[(ii)] \textbf{Fekete.} $\phi_n\ge0$, because each per-cell term
is a conditional mutual information; hence $f\ge0$ is subadditive and
\[
L^\infty(\Delta)\;:=\;\lim_{n\to\infty}\phi_n(\Delta)
\;=\;\inf_n\phi_n(\Delta)
\]
\emph{exists} in $[0,\infty)$.
\end{enumerate}
\end{namedthm}

\noindent\emph{Hygiene (both steps were demanded by R-IND-5 and are
stated, not assumed).} (a) The non-negativity $\phi_n\ge0$ is what
makes Fekete's lemma give a finite limit rather than $-\infty$; it is
immediate from the CMI form of the coordinate, but it is a hypothesis
of the lemma and is recorded as such. (b) The concatenation argument
is run on $\varepsilon$-optimal records and the limit
$\varepsilon\downarrow0$ is taken at the end, so the theorem never
needs $\phi_n$ to be \emph{attained}. (Attainment does hold, by
Theorems~\ref{thm:repr} and~\ref{thm:cvx} plus the moment box, but it
is not used.)

\begin{corollary}[The lower transfer side is a theorem]
\label{cor:lower}
$0\le\phi_n(\Delta)-L^\infty(\Delta)$ for every $n$ --- unconditionally.
Every certified finite-$n$ family value of \S\ref{sec:anchors} is
therefore an upper bound on the limit; in particular
$L^\infty(0)\le\phi_{32}(0)\le0.5647419677$.
\end{corollary}

\subsection{The reverse direction: the cross-block read $X$}
\label{sec:repair}
The matching \emph{upper} transfer bound $\phi_n-L^\infty\le c/n$ needs
the boundary charge to be controlled at the optimizers of the doubling
chain. At v0.5 that control was a bare hypothesis, because
``$\le\kappa$ per side'' is refuted as a family lemma
(Remark~\ref{rem:refute}). This revision \emph{repairs} it: a lemma
chain, using no optimizer structure whatsoever, reduces the whole
reverse direction to a single scalar.

\begin{definition}[The cross-block read]
\label{def:X}
For a split of the $n$-window into $b_1=\{1,\dots,m\}$ and
$b_2=\{m+1,\dots,n\}$ write $W^{b_i}=T^{b_i}=(V^{b_i},Y^{b_i})$,
$E=(S^{b_1},\hat Y^{b_1})$ and $E'=S^{b_2,1..\Delta}$ (empty at
$\Delta=0$). The \emph{cross-block read} of a record is
\[
X\;:=\;I\bigl(\hat Y^{b_1};\,W^{b_2}\ \big|\ W^{b_1}\bigr)\;\ge\;0 ,
\]
the information the block-1 \emph{rows} draw from block-2 sources
beyond what block-1 sources already say. $X=0$ exactly when the record
is \emph{split-causal} at $m$ (Definition~\ref{def:sc}).
\end{definition}

\begin{lemma}[L1 --- $\kappa$ is an \emph{identity}, not a bound]
\label{lem:L1}
$I(T^{b_1};T^{b_2})=\kappa$ \emph{exactly}, at every split and every
$n$. \emph{Proof.} $V$ is AR(1), hence Markov, so
$T^{b_1}\to(V_m,V_{m+1})\to T^{b_2}$ and both data-processing steps are
\emph{equalities}; the surviving term is
$I(V_m;V_{m+1})=\tfrac12\log_2\frac1{1-a^2}=\kappa$. Certificate:
$\max_m|I-\kappa|=3.8\times10^{-14}$ over $m\in\{3,4,6,8,12,16\}$, and
the spread over $m$ is $4.1\times10^{-14}$ --- the proof predicts a
constant, and no drift in $m$ is observed.\hfill$\square$
\end{lemma}

\begin{lemma}[L2 --- the cross-read bound, and it is sharp]
\label{lem:L2}
$I(E;T^{b_2})\le\kappa+X$ and $I(E';E)\le\kappa+X$.
\emph{Proof.} $S^{b_1}$ drops out of the numerator because
$U\perp(W,\hat Y)$: conditionally on $W^{b_1}$ the reference noise
carries nothing about $T^{b_2}$. What is left is
$I(W^{b_1};T^{b_2})+I(\hat Y^{b_1};W^{b_2}\mid W^{b_1})=\kappa+X$ by
Lemma~\ref{lem:L1} and Definition~\ref{def:X}; the second leg is the
same computation with $E'$ in place of $T^{b_2}$.\hfill$\square$
\emph{Sharpness (this is the point).} Over the pinned battery ---
generic $\mathcal F_0$ records, split-causal records, and the two
adversarial copy-row classes of Remark~\ref{rem:refute} --- the minimum
slack is $+5.67\times10^{-7}$, attained \emph{at the two-$V$-copy
counterexample}, i.e.\ the bound is tight exactly where the universal
claim dies. (The second leg's minimum slack is $+0.365$.)
\end{lemma}

\begin{lemma}[L3 --- the block-2 charge]
\label{lem:L3}
$D_2\;\le\;I(E;T^{b_2})-I_m$. \emph{Proof.} The interleaved chain rule
gives the exact identity
\[
D_2\;=\;I(E;T^{b_2})-I\bigl(E;T^{b_2}\mid S',\hat Y'\bigr)
-\sum_{j}I\bigl(E;S'_j\mid \mathrm{pfx}_j\bigr) ,
\]
whose first subtracted term is a mutual information ($\ge0$); the
second sum has every term $\ge0$, so it is at least its partial sum
over $j\le\Delta+1$, which telescopes (there $\mathrm{pfx}_j=S'^{\,
1..j-1}$, no record cell having been erased yet) to
$I(E;S'^{\,1..\Delta+1})\ge I_m$ by data processing
($E\supseteq S^{b_1}$). Certificates: identity residual
$\le9.3\times10^{-11}$, chain-rule collapse $6.7\times10^{-16}$,
$\min\bigl(I(E;S'^{\,1..\Delta+1})-I_m\bigr)=+1.47\times10^{-2}$ over
the battery, and
the resulting bound holds over the whole battery with minimum slack
$+0.528$. The zero-claim of Remark~\ref{rem:f0only} enters here and is
$\le2.5\times10^{-10}$ on $\mathcal F_0$.\hfill$\square$
\end{lemma}

\begin{lemma}[L4 --- the block-1 charge]
\label{lem:L4}
$D_1\le I(E';E)$, and $D_1=0$ \emph{exactly} at $\Delta=0$ (the
block-1 extension set is empty --- measured $0$ to $0.0$).
\hfill$\square$
\end{lemma}

\begin{namedthm}[H --- the repair inequality;
\emph{unconditional on $\mathcal F_0$}]
\label{thm:hrep}
For every record of $\mathcal F_0$ and every split at $m$,
\[
\boxed{\;D_1+D_2\;\le\;c(\Delta;m)+\bigl(2-\mathbf 1\{\Delta=0\}\bigr)X\;}
\qquad
c(\Delta;m)=\bigl(2-\mathbf 1\{\Delta=0\}\bigr)\kappa-I_m ,
\]
$I_m=I(S^{b_1};S'^{\,1..\Delta+1})$: add Lemmas~\ref{lem:L3}
and~\ref{lem:L4}, and bound each right-hand side by
Lemma~\ref{lem:L2}. \emph{No optimizer structure is used anywhere in
the chain} --- it is a deterministic, analytic statement about the
family. Over the pinned battery (22 cells, \emph{including both
refuting counterexamples}) there are zero violations, minimum slack
$+0.768$. The index of $c$ is the \emph{block} length; since
$I$ is monotone increasing (Remark~\ref{rem:const}) the base value is
the supremum along the dyadic chain, and numerically
$c(\Delta;8)=c(\Delta;32)$ to $2\times10^{-8}$ --- the constants of
v0.5 are unchanged.
\end{namedthm}

\begin{definition}[The split-causal sub-family]
\label{def:sc}
$\mathcal F_0^{\mathrm{sc}}(m)\subset\mathcal F_0$ is the set of
records whose block-1 cells read no block-2 source; in moment
coordinates this is the \emph{linear} condition
\[
H[b_1,\mathrm{idx}_2]\;=\;H[b_1,\mathrm{idx}_1]\,
\Sigma_{11}^{-1}\Sigma_{12},
\]
$\mathrm{idx}_i$ the source indices of block $i$ and
$\Sigma=\Sigma_W$. Equivalently $X=0$.
\end{definition}

\begin{corollary}[(H$^\ast$) is a \emph{theorem} on
$\mathcal F_0^{\mathrm{sc}}(m)$, with the same constants]
\label{cor:scthm}
On $\mathcal F_0^{\mathrm{sc}}(m)$, $X=0$ and hence
$D_1+D_2\le c(\Delta;m)$ --- exactly the inequality (H$^\ast$)
asserts, now proved, with this document's own constants. Numerically:
$X$ vanishes to $1.0\times10^{-14}$ on projected records and the
minimum slack of $D_1+D_2\le c(\Delta;m)$ over them is $+0.768$.
Moreover $\mathcal F_0^{\mathrm{sc}}(m)$ is a \emph{convex linear
section} of $\mathcal D$ (moment midpoints of split-causal records stay
split-causal: residual $1.1\times10^{-16}$, midpoint $X\le
5.1\times10^{-15}$), so the Lagrangian machinery of
Theorem~\ref{thm:cvx} restricted to the section certifies
$\phi^{\mathrm{sc}}$ two-sided exactly as it does for the diagonal
section (Remark~\ref{rem:diagsection}).
\end{corollary}

\begin{center}
\begin{tabular}{lll}
$(n,\Delta)$ & certified bracket for $\phi^{\mathrm{sc}}_n$
  & split-causality price $n\,\delta_n$ (bits)\\
\hline
$(16,0)$ & $[0.5669667555,\ 0.5670144753]$ & $0.004101$\ \ ($\le0.004165$)\\
$(24,0)$ & $[0.5655717237,\ 0.5655854048]$ & $0.004117$\ \ ($\le0.004207$)\\
$(32,0)$ & $[0.5647601440,\ 0.5648697385]$ & $0.004089$\ \ ($\le0.004382$)\\
$(16,1)$ & $[0.5414842536,\ 0.5414911968]$ & $0.010954$\ \ ($\le0.011002$)\\
$(16,2)$ & $[0.5367829566,\ 0.5367912904]$ & $0.014358$\ \ ($\le0.014408$)\\
\end{tabular}
\end{center}

\noindent Here $\delta_n=\phi^{\mathrm{sc}}_n-\phi_n$ and the
parenthesised figure is the certified (LB-to-UB) upper bound on the
price. \textbf{The price is flat in $n$} at $\approx0.0041$ bits at
$\Delta=0$ across $n=16,24,32$: split-causality costs a pure
\emph{boundary} constant, not a rate --- which is what makes it a
plausible object for a transfer argument in the first place.

\begin{lemma}[$\Pi$ --- the split-causal retraction]
\label{lem:pi}
Let $\Pi$ replace the block-1 rows of a record by their
$W^{b_1}$-conditional means plus independent noise matched to restore
$\operatorname{Cov}(\hat Y^{b_1})$. Then $\Pi$ maps $\mathcal F_0$ onto
$\mathcal F_0^{\mathrm{sc}}(m)$ and
\begin{enumerate}
\item[(i)] preserves the per-cell \emph{distortion} exactly
(measured $1.1\times10^{-16}$);
\item[(ii)] preserves the entire block-1 joint law
$(W^{b_1},S^{b_1},\hat Y^{b_1})$ exactly ($6.7\times10^{-16}$), hence
every own-window block-1 term;
\item[(iii)] kills the cross-read, $X(\Pi R)\le2.1\times10^{-14}$;
\item[(iv)] \emph{strictly improves} generic records --- over the
pinned generic battery ($n\in\{12,16,20\}$, three styles,
$\Delta\in\{0,2\}$) the gain $n\bigl(L_a(\Pi R)-L_a(R)\bigr)$ lies in
$[-2.74,-0.31]$ bits (prover's own battery: $[-3.11,-0.37]$);
\item[(v)] costs only $\approx5.0\times10^{-3}$ bits \emph{at
optimizers} ($0.005052/0.005044/0.005039/0.005035$ at
$m=8/12/16/24$, $\Delta=0$).
\end{enumerate}
$\Pi$ is what makes $\mathcal F_0^{\mathrm{sc}}$ a \emph{retract} and
not merely a subset, and (i)--(iii) are what make the certified
$\phi^{\mathrm{sc}}$ brackets above meaningful. \textbf{What $\Pi$
does \emph{not} do is bound $X$}: at an optimizer its price is
$\ge0$ automatically, so ``price $\ge0$'' is consistent with any $X$
whatever --- see Remark~\ref{rem:repair}(b), where that route is
closed. The measured price/$X$ ratio is $0.7960$ at every $n$ and
every $\Delta$.
\end{lemma}

\begin{hypothesis}[(H$^{\ast\ast}$) --- the cross-block read at
dyadic-chain optimizers; hard, load-bearing]
\label{hyp:hstarstar}
Fix $\Delta$ and a base window length $n$. At every optimizer of the
dyadic chain $n\cdot2^k$, $k=0,1,2,\dots$, the cross-block read of the
split into halves satisfies $X\le\bar X$.
\end{hypothesis}

\begin{corollary}[(H$^\ast$) --- what v0.5 assumed --- is now a
\emph{consequence}]
\label{cor:hstar}
Under (H$^{\ast\ast}$), Theorem~\ref{thm:hrep} gives
\[
D_1+D_2\;\le\;c(\Delta;n)+\bigl(2-\mathbf 1\{\Delta=0\}\bigr)\bar X
\;=:\;c^\ast(\Delta;n) ,
\]
which is (H$^\ast$) with the constant enlarged by the repair term.
(H$^{\ast\ast}$) is \emph{strictly weaker and better margined}: it
constrains \emph{one scalar} rather than a sum of $2m$ conditional
mutual informations, the measured value at the $\Delta=0$ optimizers is
$X=6.35\times10^{-3}$ against the largest admissible
$\bar X=32\bigl(\mathrm{LB}(32,0)-0.5479448\bigr)-c(0)=0.116930$ ---
$18\times$ headroom, against (H$^\ast$)'s $0.076971$ inside $0.4202858$
($5.5\times$) --- it is verified \emph{monotone decreasing} in $m$
(Remark~\ref{rem:mtable}), and it is constructively bounded at any $m$
by $\Pi$ (Lemma~\ref{lem:pi}).
\end{corollary}

\begin{namedthm}[The split-causal transfer step; \emph{unconditional}]
\label{thm:scstep}
For every $m$ and every $\Delta$,
\[
\phi^{\mathrm{sc}}_{2m}(\Delta)\;\le\;\phi_m(\Delta) :
\]
concatenate two independent $\varepsilon$-optimal $m$-window records.
The result is feasible on the $2m$-window (the distortion constraint is
the cell average) and is \emph{split-causal at the midpoint} --- its
block-1 rows read no block-2 source whatever --- and its value is at
most $\phi_m+\varepsilon$ by the per-cell inequality of
Theorem~\ref{thm:transfer}(i). Certified form (upper end of
$\phi^{\mathrm{sc}}_{2m}$ against lower end of $\phi_m$): the
inequality holds with certified slack
$3.8/2.4/1.9\times10^{-3}$ at $(2m,m)=(16,8)/(24,12)/(32,16)$.
Applying Corollary~\ref{cor:scthm} to the $2m$-window
$\mathcal F_0^{\mathrm{sc}}$-optimizer gives
$\phi_m-\phi^{\mathrm{sc}}_{2m}\le c(\Delta;m)/2m$, whence
\[
\phi_m(\Delta)-\phi_{2m}(\Delta)\;\le\;
\frac{c(\Delta;m)+p_{2m}}{2m},
\qquad
p_n\;:=\;n\bigl(\phi^{\mathrm{sc}}_n-\phi_n\bigr)
\]
the \emph{split-causality value gap} --- the same quantity as the
price column of the $\phi^{\mathrm{sc}}$ table above.
\end{namedthm}

\begin{hypothesis}[(H$^{\ast\ast\ast}$) --- the split-causality value
gap along the chain]
\label{hyp:hstar3}
$p_m\le\bar p$ for every $m$ in the dyadic chain.
\end{hypothesis}

\begin{remark}[Three hypotheses, and why (H$^{\ast\ast\ast}$) is the
one to prefer]
\label{rem:hchain}
Each of (H$^\ast$), (H$^{\ast\ast}$), (H$^{\ast\ast\ast}$) suffices for
Theorem~\ref{thm:transfer2}, and they are ordered by how little they
assume and --- more importantly --- by how \emph{checkable} they are.
(H$^{\ast\ast}$) implies (H$^\ast$) with an enlarged constant, by
Theorem~\ref{thm:hrep} (Corollary~\ref{cor:hstar}); it replaces a bound
on a sum of $2m$ conditional mutual informations by a bound on one
scalar. (H$^{\ast\ast\ast}$) replaces that scalar by another,
$p_m$, which reaches the same conclusion through
Theorem~\ref{thm:scstep} rather than through
Theorem~\ref{thm:hrep} --- and \textbf{$p_m$ is a gap between two
convex programs, so Theorem~\ref{thm:cvx} certifies it two-sided},
whereas $X$ is only ever \emph{measured} at a computed optimizer. That
is the decisive difference: a hypothesis about $p$ is falsifiable by
certificate at any $m$ one is willing to pay for. Both bars are the
same number, $\bar p,\ \bar X<32\bigl(\mathrm{LB}(32,0)-0.5479448
\bigr)-c(0)=0.116930$; the certified instances --- each the
LB-to-UB upper end, so each an honest \emph{bound} on $p_n$ rather
than an optimizer reading --- are $p_{16}\le0.004165$,
$p_{24}\le0.004207$, $p_{32}\le0.004382$: flat in $n$, and
$27\times$ inside the bar, against $X$'s $18\times$.
\end{remark}

\begin{remark}[The ``$\le\kappa$ per side'' refutation, \emph{refined}
--- what it refutes is split-causality]
\label{rem:refute}
The natural lemma --- that each side's boundary charge is at most the
AR(1) boundary information $\kappa$, uniformly over the family --- is
\emph{false}, and its sharpened form $D_i\le\kappa-I_m$ is false too.
Explicit $D$-feasible records in $\mathcal F_0$ (no $U$-coupling
whatsoever; the whole budget respected) spend block-1 rows on copies of
block-2 cells:
\begin{itemize}
\item two $V$-copy rows give $D_2=0.4563>0.4203=\kappa-I_8$: the
\emph{sharpened} constant is exceeded by a mild, feasible
perturbation;
\item three $Y$-copy rows give $D_2=18.93$ bits $\approx26\kappa$ at
distortion $0.2986$: the charge is not bounded by $\kappa$ at all.
\end{itemize}
Nothing stops the construction from being repeated: the worst case over
the family is $\Theta(m)$ in the block length. \textbf{No universal
constant exists} --- that statement stands unchanged.

\emph{What has changed is the diagnosis.} These are precisely the
records with enormous cross-block read: the harness's own witnesses of
the same class measure $X=18.46$ and $X=29.27$ bits (against
$6.3\times10^{-3}$ at the optimizers), and at the two-$V$-copy witness
the bound of Lemma~\ref{lem:L2} is \emph{tight to
$5.7\times10^{-7}$}. The refutation is exactly the failure of
split-causality, and Theorem~\ref{thm:hrep} absorbs it into the term
$(2-\mathbf 1\{\Delta=0\})X$ rather than being defeated by it. Hence
the v0.5 blanket rule (``this document never prints $\le\kappa$ per
side as an $\mathcal F_0$ lemma'') is \emph{refined}, not repealed:
$D_1+D_2\le c(\Delta;m)$ is a \emph{theorem} on
$\mathcal F_0^{\mathrm{sc}}(m)$ (Corollary~\ref{cor:scthm}) and is
\emph{false} on $\mathcal F_0$; neither form may be printed without its
scope. What no longer follows is the v0.5 sentence ``no proof of the
reverse direction can avoid using optimizer structure'': the
\emph{inequality} of Theorem~\ref{thm:hrep} avoids it entirely; what
still needs optimizer structure is a bound on $X$
(Remark~\ref{rem:repair}).

The two constants above are the verifier's witnesses; they are not in
the repository, so the harness pins its own analytic records of the
same class and measures $D_2=0.7749$ for two $V$-copy rows (distortion
$0.2706$, $L_a(0)=4.219$) and $D_2=13.94$ bits $\approx19\kappa$ for
three $Y$-copy rows (distortion $0.2754$, $L_a(0)=4.740$). Same
refutation, independent witnesses; the refutation itself is a gated
fact of the harness, not a remark. Note the $L_a$ values: these records
are nowhere near optimal --- more than seven times $\phi_{16}=0.5668$
--- which is why they never threatened the transfer bound in practice,
only its proof.
\end{remark}

\begin{remark}[The monotone $m\times\Delta$ table (the
``$m\in\{8,16\}$ only'' caveat, retired)]
\label{rem:mtable}
v0.5 recorded the optimizer verifications at $m\in\{8,16\}$ only and
flagged the gap explicitly. The table is now
$m\in\{8,12,16,24\}$ at $\Delta=0$ and $m\in\{8,12,16\}$ at
$\Delta=1,2$ (cold-start certified full-space optimizers at $n=2m$),
and \emph{every} quantity of the repair is monotone
\emph{decreasing} in $m$:
\begin{center}
\begin{tabular}{lllll}
$\Delta$ & $m$ & $D_1+D_2$ & $I(E;T^{b_2})$ & $X$\\
\hline
$0$ & $8/12/16/24$ & $0.076971/0.076934/0.076915/0.076896$
 & $0.543069/0.543059/0.543054/0.543049$
 & $6.3468/6.3360/6.3307/6.3254\ (\times10^{-3})$\\
$1$ & $8/12/16$ & $0.090533/0.090487/0.090463$
 & $0.569542/0.569504/0.569485$
 & $1.8013/1.7985/1.7971\ (\times10^{-2})$\\
$2$ & $8/12/16$ & $0.099029/0.098964/0.098931$
 & $0.579319/0.579256/0.579225$
 & $2.3894/2.3849/2.3826\ (\times10^{-2})$\\
\end{tabular}
\end{center}
Monotonicity is what the caveat was about: \emph{the base value
dominates the whole chain}, so a hypothesis verified at the base $m$
is not being extrapolated upward along the doubling chain.
\emph{Scope of the netting:} the $\Delta=0$ row is reproduced
end-to-end and gated (strict decrease at every step, each step
$\ge10^{-6}$ against a smallest measured step of
$4.96\times10^{-6}$); the $\Delta=1,2$ rows are the prover's, carried
on its authority and not re-netted, except that the $m=8$ optimizers
at $\Delta=1,2$ are recomputed for Theorem~\ref{thm:K}. Two-point
$1/m$ limits at $\Delta=0$: $X_\infty\approx6.31\times10^{-3}$,
$(D_1+D_2)_\infty\approx0.076859$,
$I(E;T^{b_2})_\infty\approx0.54304$; the $\Pi$-price of
Lemma~\ref{lem:pi}(v) is monotone too, limit $\approx5.03\times10^{-3}$.

\emph{Two precision notes on the values v0.5 recorded at $m=16$.}
(a) The recorded $D_2=0.076926$ and $I(E;T^{b_2})=0.543057$ came from a
\emph{warm-polish} endpoint sitting at $d-D=+6.4\times10^{-5}$, i.e.\
\emph{distortion-infeasible}; the feasible cold-start certified values
are $0.0769149$ and $0.543054$. (b) The difference, $1.1\times10^{-5}$
resp.\ $3\times10^{-6}$, is more than two orders inside the $5\times
10^{-3}$ / $2\times10^{-2}$ reproduction bands the 080 harness gates
these two quantities at --- so this is a \emph{precision note, not a
bar change}, and the sealed 080 verdicts are untouched. The table above
quotes the feasible cold-start values throughout.
\end{remark}

\begin{namedthm}[T$^\ast$ v2 --- boundary-charged reverse direction and
the two-sided bound, in \emph{two} forms]
\label{thm:transfer2}
Fix $\Delta$ and a base $n$, and write
$c^\ast(\Delta;n)=c(\Delta;n)+\bigl(2-\mathbf 1\{\Delta=0\}\bigr)
\bar X$.
\begin{enumerate}
\item[(ii-a)] \emph{Family form, conditional on
Hypothesis~\ref{hyp:hstarstar} along the chain $n\cdot2^k$:}
$\displaystyle
\phi_m(\Delta)-\phi_{2m}(\Delta)\;\le\;\frac{c^\ast(\Delta;n)}{2m}$
for every $m=n\cdot2^k$, and, summing the geometric telescope,
$\displaystyle 0\;\le\;\phi_n(\Delta)-L^\infty(\Delta)\;\le\;
\frac{c^\ast(\Delta;n)}{n}$, the lower half being
Corollary~\ref{cor:lower} and \emph{unconditional}.
\item[(ii-b)] \emph{Split-causal form, \emph{unconditional}:} the same
two inequalities hold for the sub-family value
$\phi^{\mathrm{sc}}_n(\Delta)=\inf_{\mathcal F_0^{\mathrm{sc}}}L_a$
with the constant $c(\Delta;n)$ itself and \emph{no hypothesis at
all}, because on $\mathcal F_0^{\mathrm{sc}}$ the charge bound is
Corollary~\ref{cor:scthm}.
\end{enumerate}
\emph{Scope of (ii-b).} It bounds
$\phi^{\mathrm{sc}}_n-L^{\mathrm{sc},\infty}$, not
$\phi_n-L^\infty$; the two coincide iff the split-causality price
$\delta_n=\phi^{\mathrm{sc}}_n-\phi_n$ is $o(1/n)$ in the relevant
sense. Measured, $n\delta_n$ is \emph{flat} at $\approx0.0041$ bits
over $n=16,24,32$ at $\Delta=0$ (so $\delta_n=\Theta(1/n)$, of the same
order as the transfer term itself, not smaller). That is evidence, not
a proof, and (ii-b) is therefore \emph{not} a route to the family
statement --- it is the unconditional theorem about the sub-family, and
the honest statement of what the repair buys without (H$^{\ast\ast}$).

\emph{Exact superadditivity is refuted}, so (ii) cannot be sharpened to
$c=0$: the certified anchors are strictly decreasing,
$\phi_8>\phi_{12}>\phi_{16}>\phi_{24}>\phi_{32}$, with the
certified transfer anchors
$\phi_8(0)\in[0.570790938,\,0.5707933842]$ and
$\phi_{12}(0)\in[0.568030932,\,0.5681028134]$. The bound is not tight
either: the measured drift $n\bigl(\phi_n-L^\infty\bigr)\approx0.0645$
at $n=32$ sits $6.5\times$ inside $c(0)=0.4203$.
\end{namedthm}

\begin{remark}[The per-block feasibility step of (ii) --- a separate
hole, and its repair]
\label{rem:feas}
The restriction step behind (ii) needs each half of the
$2m$-optimizer to be a \emph{feasible} $m$-window record, i.e.\
$\mathrm{own}_i\ge m\,\phi_m(D)$, and v0.5 tacitly assumed each block
distortion is $\le D$. \textbf{It is not.} The $2m$-optimizer's block
distortions \emph{straddle} $D$ --- $d_1=0.299486$, $d_2=0.300514$ at
$n=16$; $0.299658/0.300342$ at $n=24$; $0.299744/0.300256$ at $n=32$
--- with the \emph{mean} exactly $D$ (to $5\times10^{-8}$). The correct
step is
\[
\mathrm{own}_1+\mathrm{own}_2\;\ge\;
m\bigl[\phi_m(d_1)+\phi_m(d_2)\bigr]\;\ge\;2m\,\phi_m(D),
\]
the second inequality by \emph{convexity of the value function}
$D\mapsto\phi_m(D)$, which Theorem~\ref{thm:cvx} supplies for free
(jointly convex objective, affine distortion constraint $\Rightarrow$
convex optimal value in the right-hand side). Certified on the pinned
grid $D\in\{0.26,0.28,0.30,0.32,0.34\}$ at $m=8$, $\Delta=0$: at each
interior node the certified \emph{upper} end sits below the chord of
the two certified \emph{lower} ends, slack
$1.90/1.64/1.47\times10^{-3}$. Ignoring the straddle would have cost
$m\,\mu\,(d_2-D)=8\times2.2206\times5.14\times10^{-4}
=9.13\times10^{-3}$ bits at $n=16$ ($\mu$ the distortion multiplier at
the optimizer) --- \emph{larger than the entire repair term} $\bar X$,
so the hole was not cosmetic. It is independent of (H$^\ast$) and of
everything else in this section.
\end{remark}

\begin{remark}[The constants, and their $n$-uniformity]
\label{rem:const}
$I_n=I(S^{b_1};S'^{\,1..\Delta+1})$ is \emph{monotone increasing in
$n$}, so $c(\Delta;n)=(2-\mathbf 1\{\Delta=0\})\kappa-I_n$ is the
supremum of the per-step constants along the whole dyadic chain issuing
from base $n$ --- which is precisely what the telescope in
Theorem~\ref{thm:transfer2}(iii) requires, and why the constant must be
built from $I_n$ and not from $I_\infty$. Numerically $I_n$ is
converged to $I_{32}$ by $n=16$ (to $<10^{-12}$), and
\[
I_{32}=0.3166798\,/\,0.3482018\,/\,0.3520686
\quad\text{at }\Delta=0/1/2,
\]
each equal to $I_\infty$ to below $10^{-15}$. Hence
\[
c(0)\le0.42029,\qquad \boxed{c(1)\le1.1258},\qquad c(2)\le1.1219 .
\]
\emph{Erratum:} the transfer note recorded $c(1)\le1.1257$; the
computed value is $1.1257294$, so $1.1257$ is false in the fifth
decimal. Use $1.1258$.

\emph{The repaired constants (v0.6).} With
$c^\ast(\Delta;n)=c(\Delta;n)+(2-\mathbf 1\{\Delta=0\})\bar X$ and
$\bar X$ at the base-$m$ measured values of Remark~\ref{rem:mtable},
$c^\ast(0)=0.4266326$, $c^\ast(1)=1.161755$, $c^\ast(2)=1.169651$ ---
the $\Delta=0$ constant grows by $1.5\%$, the others by $3$--$4\%$,
because $X$ itself grows with $\Delta$ and enters doubled there. Only
$c^\ast(0)$ is used in Corollary~\ref{cor:plateau}; $\Delta=1,2$ remain
$O(1/n)$-model-conditional for the reason given there, and the repair
does not change that.
\end{remark}

\begin{remark}[The zero-claim and the marginalization step are
$\mathcal F_0$-only]
\label{rem:f0only}
Two steps carry the whole cancellation. (a) The \emph{zero-claim}
$I\bigl(E;S'_j\mid \mathrm{pfx}_j,T^{b_2}\bigr)=0$ makes the interleaved
chain rule cancel the record--reference coupling exactly (residual
$\le2\times10^{-14}$ over pinned $\mathcal F_0$ records). (b) The
\emph{marginalization} step replaces the block-2 sub-record by an
effective $8$-window record with the same law (deviation
$2\times10^{-16}$; the effective noise is independent of both
$W^{b_2}$ and $U^{b_2}$).
\textbf{Both are tagged $\mathcal F_0$-only.} They fail outside the
family, and the failure is not a technicality: with block-1 rows
coupled to $U^{b_2}$ the zero-claim is violated by $0.2223$ bits, and
with the coupling confined to \emph{block-2 rows only} it is still
violated by $0.0950$ bits --- in the latter case the violation is
induced purely through the conditioning, with no block-1 record
touching $U$ at all. So $U$-independence is load-bearing here in the
same way it is load-bearing in Theorem~\ref{thm:repr}
(Remark~\ref{rem:ulegs}), and for a third time.
\end{remark}

\begin{corollary}[The $\Delta=0$ plateau over the causal-spectral
allocation --- recomputed with the repaired constant]
\label{cor:plateau}
Assume Hypothesis~\ref{hyp:hstarstar} at $\Delta=0$ with base $n=32$,
with $\bar X$ taken at the measured value $X=6.346782\times10^{-3}$.
With the sealed 079 certified lower end
$\mathrm{LB}(32,0)=0.5647327869$ (harness
\texttt{go14\_convexity.py}, entry \texttt{s5\_32\_0}) and
$c^\ast(0)=c(0)+X=0.4202858+0.0063468=0.4266326$,
\[
L^\infty(0)\;\ge\;0.5647327869-\frac{0.4266326}{32}
\;=\;0.5514005\;>\;0.5479448,
\]
the right-hand side being the causal-spectral allocation at $\Delta=0$
(Remark~\ref{rem:refcon}). The margin is $+3.456\times10^{-3}$:
$203\times$ the $1.7\times10^{-5}$ bracket width quoted for $(32,0)$ in
the anchor table of \S\ref{sec:anchors}, and $376\times$ the
$9.18\times10^{-6}$ width the sealed harness actually achieved.
\emph{This is a statement about the limit itself, not about the
$O(1/n)$ finite-window model.} (At v0.5, under (H$^\ast$), the same
arithmetic gave $0.5515989$ at $215\times$/$398\times$; the repair
costs $1.98\times10^{-4}$ of margin and buys a hypothesis with
$18\times$ headroom instead of $5.5\times$.)

\emph{Under the decay route instead.} Replacing (H$^{\ast\ast}$) by the
two model-scale hypotheses (D) and (F) of Remark~\ref{rem:repair} gives
the proved envelope $\bar X=0.065038$, hence $c^\ast(0)=0.485324$ and
\[
L^\infty(0)\;\ge\;0.5647327869-\frac{0.485324}{32}\;=\;0.5495664
\;>\;0.5479448 ,
\]
margin $+1.622\times10^{-3}$ ($95\times$ the quoted bracket,
$177\times$ the sealed width). Which hypothesis buys which margin is
therefore explicit: (H$^{\ast\ast}$) --- one measured scalar ---
buys $3.456\times10^{-3}$; (D)$+$(F) --- two structural, $m$-uniform
statements with a \emph{proved} conversion --- buys
$1.622\times10^{-3}$. Neither is discharged.

\emph{The margin is not slack in the constant, and the base matters.}
Running the same arithmetic from base $n=24$ gives
$0.5654101216-0.4266326/24=0.5476338<0.5479448$: it \emph{fails} by
$3.11\times10^{-4}$ (and by $2.757\times10^{-3}$ under (D)$+$(F)) ---
so the enlarged constant makes the no-slack claim \emph{stronger} than
at v0.5, where the base-24 shortfall was $4.7\times10^{-5}$. The
hypothesis is genuinely load-bearing, and so is the choice of base
window. $\Delta=1$ and $\Delta=2$ remain
$O(1/n)$-model-conditional: their constants are $2.7\times$ larger and
their plateaus $3.6$/$18\times$ smaller, so a certified anchor at
$n\approx260$ resp.\ $n\approx1320$, or a sharper constant, would be
needed.
\end{corollary}

\begin{remark}[The reduction, executed --- and the obstruction, named]
\label{rem:repair}
v0.5 recorded a \emph{repair path}. It has since been walked, and it
ends at a named theorem rather than at a vague hope. Three things are
now on record.

\emph{(1) The conversion is proved and is $m$-uniform by
construction.} Exactly,
$X=\tfrac12\log_2\det\bigl(I+N_{11}^{-1}A^{12}\Sigma_{2\mid1}
(A^{12})^{\!\top}\bigr)$ with $N$ the record noise covariance and
$A^{12}$ the cross-block kernel. Under
\begin{itemize}
\item[(D)] \emph{kernel-decay envelopes}: $|A_v[t,s]|\le
C_v\lambda^{|t-s|}$, $|A_y[t,s]|\le C_y\lambda^{|t-s|}$ with
$\lambda=\lambda_s=0.354554$, $C_v=0.1980$, $C_y=0.1631$
(fitted, and identical to three digits at $n=24/32/48$), and
\item[(F)] a uniform noise floor $\lambda_{\min}(N)\ge\nu=0.1495$,
\end{itemize}
one gets a closed-form bound on $X$ whose sums run over infinite
half-lines and which is therefore \emph{independent of $m$}; the
conditional covariance $\Sigma_{2\mid1}$ is used exactly rather than
enveloped (its $-a^{j+k+2}$ boundary term is worth $2.6\times$ and is
what makes the bound clear the bar). Audited link by link at real
optimizers the total looseness is $10.3\times$
($\log\det\to\operatorname{tr}$ $1.004\times$, $N_{11}\succeq\nu I$
$1.22\times$, the enveloped Schur step $8.4\times$). \textbf{Result:
$\bar X=0.065038$ against the admissible $0.116930$ --- it clears with
$1.80\times$ margin, uniformly in $m$} (free-$\lambda$ optimisation
gives $0.0648$; insensitive to the envelope-fit window). So
(H$^{\ast\ast}$) is no longer a bare hypothesis but a
\emph{consequence} of (D) and (F), at the arithmetic cost recorded in
Corollary~\ref{cor:plateau}.

\emph{(2) The spectral premise (D) is not the obstruction --- it
holds numerically, uniformly in $m$.} At $\Delta=0$, $m=8/12/16/24$:
$\theta$ moves only over $3.0785$--$3.0835$, $\operatorname{cond}N$
over $1.4453$--$1.4492$, $\operatorname{cond}M_Q$ over
$2.197$--$2.217$, $\lambda_{\min}(J)=0.2686$ flat, the record pivots
$\sigma_t\in[0.3753,0.4588]$, and $\|A^{12}\|_F$ is \emph{flat in $m$}
($5.4823\to5.4722\times10^{-2}$; $\|A^{12}\|_F/\|A^{11}\|_F=3.44\%$).
The measured kernel is the same Toeplitz kernel plus non-moving
boundary corrections at every $n$, and its per-lag ratios
\emph{decrease} with lag ($0.284\to0.200$ for $A_y$,
$0.377\to0.279$ for $A_v$), so a single geometric envelope at
$\lambda_s$ is conservative.

\emph{(3) The obstruction is \textbf{circularity of the operator}, and
it is a genuine research step.} By Theorem~\ref{thm:K} the only
operator whose inverse produces the optimizer is
$N=(\theta I+V\Sigma_s^{-1}V^{\!\top})^{-1}$, and
$M:=V\Sigma_s^{-1}V^{\!\top}=N^{-1}-\theta I$ \emph{exactly}: ``$M$
decays exponentially'' and ``$N$ decays exponentially'' are the same
statement, so a Demko/Jaffard inverse-closedness theorem applied here
\emph{returns its own hypothesis}. The model's own decay ($\Sigma_S$,
$\Sigma_W$, $\Sigma_{W\mid S}$ all exactly $a^{|t-s|}$) enters the KKT
system only through the affine data, never through the operator being
inverted. What would inject it is a \emph{fixed-point/bootstrap}
argument in a Jaffard (or Gr\"ochenig--Klotz weighted) Banach algebra:
show that the KKT map sends a ball
$\{\|A\|_{\mathcal A_\lambda}\le R\}$ into itself and contracts. A
second, independent piece is missing too: (F), the uniform lower bound
$\lambda_{\min}(N)\ge\nu$. Theorem~\ref{thm:K} gives the \emph{upper}
side ($N\preceq I/\theta$, proved) but not the lower; every attempted
route closes on itself ($\lambda_{\max}(M)\le1/\lambda_{\min}(M_Q)$
with $M_Q\succeq N$ gives only $1/\nu\le\theta+1/\nu$), and coercivity
fails because a record that does not use a direction pays nothing
there.

\emph{Routes closed, so that they are not re-attempted.} (a) KKT
zeroing cannot work: since $\Sigma_W^{-1}\Sigma_{WY}=[0;I]$, the
distortion gradient \emph{vanishes identically} on the cross-read
coordinates (Theorem~\ref{thm:K}), so stationarity there carries no
multiplier --- the cross-read is a distortion-\emph{free} direction on
which the objective is stationary, and the optimizer chooses $X>0$
because it strictly lowers $L_a$. KKT yields an equation, not a bound.
(b) The $\Pi$-retraction cannot bound $X$ either: at an optimizer the
$\Pi$-price is $\ge0$ automatically, and the measured price/$X$ ratio
is $0.7960$ at every $n$ and every $\Delta$, so ``price $\ge0$'' is
consistent with \emph{any} $X$. (c) $m$-induction restates the same
missing estimate: $X$ is monotone in $m$
(Remark~\ref{rem:mtable}) but the step compares the $m$- and
$2m$-optimizers. \emph{The next route on record is different in kind}:
solve the stationary system (K1)--(K3) in the spectral domain with the
GO-13 machinery, where it becomes per-frequency scalar equations ---
rational symbols with poles off the unit circle would give Toeplitz
coefficients decaying at exactly the modulus of the nearest pole, which
the numerics place at $\lambda_s=a(1-K)=0.354554$, GO-14's own netted
constant --- and then transfer stationary $\to$ finite-window by a
stability estimate from the strong convexity of
Theorem~\ref{thm:cvx}.

\emph{What this does \emph{not} do.} Corollary~\ref{cor:plateau}
cannot drop its hypothesis. It trades (H$^{\ast\ast}$) for (D)$+$(F);
the plateau survives either way, at $+3.456\times10^{-3}$ resp.\
$+1.622\times10^{-3}$, and the base-24 control fails either way.
\end{remark}

\subsection{The within-class ladder, certified (the
Remark~\ref{rem:diagsection} upgrade, executed)}
\label{sec:ladder}
The diagonal class is a convex \emph{linear} section of $\mathcal D$
(Remark~\ref{rem:diagsection}), so the Lagrangian machinery of
Theorem~\ref{thm:cvx} restricted to the section --- with a section
moment box $R_{\mathrm{box}}=73$/$95$/$136$ at $n=24/32/48$, verified
attained and hence tight --- certifies the ladder two-sided. All nine
ladder cells and the three block values are now certified
\emph{class} values (per cell the tighter of the two independent
end-to-end computations):

\begin{center}
\begin{tabular}{lll}
$(n,\Delta)$ & certified within-class bracket & width\\
\hline
$(24,4)$ & $[0.5333160113,\ 0.5333170450]$ & $1.0\times10^{-6}$\\
$(24,5)$ & $[0.5332969026,\ 0.5332971560]$ & $2.5\times10^{-7}$\\
$(24,6)$ & $[0.5332945088,\ 0.5332947270]$ & $2.2\times10^{-7}$\\
$(32,4)$ & $[0.5324808853,\ 0.5324814627]$ & $5.8\times10^{-7}$\\
$(32,5)$ & $[0.5324606031,\ 0.5324610143]$ & $4.1\times10^{-7}$\\
$(32,6)$ & $[0.5324578912,\ 0.5324584264]$ & $5.4\times10^{-7}$\\
$(48,4)$ & $[0.5316458819,\ 0.5316462767]$ & $3.9\times10^{-7}$\\
$(48,5)$ & $[0.5316244073,\ 0.5316250208]$ & $6.1\times10^{-7}$\\
$(48,6)$ & $[0.5316217661,\ 0.5316233575]$ & $1.6\times10^{-6}$\\
$\mathrm{block}_{24}$ & $[0.5332936495,\ 0.5332949450]$
  & $1.3\times10^{-6}$\\
$\mathrm{block}_{32}$ & $[0.5324577110,\ 0.5324581027]$
  & $3.9\times10^{-7}$\\
$\mathrm{block}_{48}$ & $[0.5316214213,\ 0.5316220475]$
  & $6.3\times10^{-7}$\\
\end{tabular}
\end{center}

\noindent The $(48,6)$ cell is the loosest (its feasible-projection
correction is $1.1\times10^{-6}$, still inside the bracket), which is
the same cell the raw-anchor demotion of Remark~\ref{rem:anchors}
concerns.

\medskip\noindent\textbf{What this does \emph{not} certify.} The excess
$E_\infty(\Delta)$ is a \emph{two-point $1/n$ extrapolation} of these
class values against the block values, so the re-assembled
$E_\infty$ intervals are \emph{$1/n$-model-conditional, not
certificates}, and are tagged as such:
$E_\infty(4)\in[2.50,2.75]\times10^{-5}$, confirming the
$2.67\times10^{-5}$ of Conjecture~\ref{conj:proc} in bar;
$E_\infty(5)\in[1.7,4.8]\times10^{-6}$, sign resolved positive by the
tighter re-assembly and \emph{unresolved} by the looser one; and
$E_\infty(6)\in[-3.2,1.5]\times10^{-6}$, whose bracket \emph{contains
zero} at every achievable width. The last of these independently
re-ratifies the R-IND-5 demotion of $E_\infty(6)$ to the raw anchor
$L_a^{\mathrm{diag}}(48,6)=0.5316222730$
(Remark~\ref{rem:anchors}). No claim is made here about how many
optimizer endpoints were true class minima; the transfer note's
``$4/9$ endpoints'' remark is \emph{withdrawn} as unverifiable from the
record.

\begin{question}[Open]
\label{q:open}
(1) \emph{A uniform-in-$m$ decay bound on the optimizer's cross-block
kernel} --- the residue of the old $n$-monotonicity face after
Theorem~\ref{thm:transfer} closed its subadditivity half
unconditionally and Theorem~\ref{thm:hrep} reduced the rest to the
single scalar $X$. The open problem is now \emph{named}, and it is not
``bound $I(E;T^{b_2})$'': (i) no \emph{uniform family} bound can exist
(Remark~\ref{rem:refute}); (ii) no KKT/exchange argument can force
$X=0$, because the distortion gradient vanishes identically on the
cross-read coordinates (Theorem~\ref{thm:K}) --- the cross-read is a
distortion-\emph{free} direction on which the objective is stationary
and the optimizer chooses $X>0$ because it strictly lowers $L_a$, so
stationarity yields an equation, not a bound; (iii) what is missing is
an \emph{optimizer decay estimate}: measured,
$\|A^{12}\|_F/\|A^{11}\|_F=3.44\%$ with the cross-block kernel decaying
geometrically at per-lag ratio $\approx0.28$--$0.35$, bracketing the
smoothing pole $\lambda_s=0.354554$, and $\|A^{12}\|_F$ is
\emph{flat in $m$}. A Demko/Jaffard-type exponential off-diagonal decay
bound for inverses of banded positive-definite operators, applied to
the KKT system of Theorem~\ref{thm:K}, would give
$\|A^{12}\|\le C\lambda^{\bullet}$ uniformly in $m$ and close the gap
--- a standard \emph{type} of theorem, much softer than
(H$^{\ast\ast}$). The conversion from such a bound to $\bar X$ is
already proved (Remark~\ref{rem:repair}); the two pieces still missing
are the fixed-point/bootstrap step that breaks the circularity
$M=N^{-1}-\theta I$, and the uniform lower bound
$\lambda_{\min}(N)\ge\nu$. Alternatively (H$^{\ast\ast\ast}$)
(Remark~\ref{rem:hchain}) may be attacked directly, since the
split-causality value gap is two-sided certifiable. Separately,
certified anchors at $n\approx260$/$1320$ would carry the
$\Delta=1,2$ plateaus without any of the three hypotheses; and the
$n$-uniform dual certificate of the GO-13 spectral machinery would
bypass the transfer route altogether;
(2) per-cell convexity (Remark~\ref{rem:percell}) --- the
within-class certification of the diagonal ladder
(Remark~\ref{rem:diagsection}) is no longer open, it is executed in
\S\ref{sec:ladder};
(3) the $U$-coupled coordinate (Hypothesis~\ref{hyp:uind} dropped):
a different, far cheaper object, definitionally open; (4) the reset
protocol realizing definition (a); (5) identification of the
closure constant beyond $\Delta\approx6$, and the innovations
characterization of the limit (vector NRDF closed forms were
refuted in general --- Stavrou--Tanaka--Tatikonda --- so only
scalar forms are expected cheaply).
\end{question}

House rules: the novelty gate is \emph{cleared} for the four
objects swept (arXiv + DBLP + Crossref; S2 residual is
belt-and-suspenders only), conditional on the five must-cite
adjacents (Berta et al.; Boyd--Mandal--Crutchfield;
Venkat--Weissman--Carmon--Shamai; Sandberg et al.; Anderson 2018)
--- headline novelty language may be promoted only with the five
citations in force, and nothing above claims it yet. R-IND-5 passes
1 (+closure), the Conjecture-v0.2 pass, the Theorems-R+C pass
(seven mandatory restatements) and the $n$-transfer pass (six
mandatory restatements --- the (H$^\ast$) demotion, the
$\kappa$-per-side refutation, the constants, the
$\mathcal F_0$-only tags, the corollary arithmetic and the
ladder tagging --- all in force in this revision) are
on record (\texttt{go14-causal-erasure-PROBE.md}); seal
GO-P-2026-076 governs the theorem's numerical faces; the
process-limit faces of \S\ref{sec:proc} await prereg 078 (harness
\texttt{experiments/go14\_process\_limit.py}); the convex-program
faces of \S\ref{sec:convex} await prereg 079 (harness
\texttt{experiments/go14\_convexity.py}: representation identity
incl.\ non-staircase schedules, PSD/closed-form/Sylvester
certificates, $4{,}200$-pair Jensen sweep, tangent plane, cold-start
certified brackets at $(16,0)/(16,1)/(24,0)/(32,0)$ +
$\mathrm{block}_{16}$, sandwich-margin window, and the two-leg /
non-conflation $U$-gates); the transfer faces of \S\ref{sec:transfer}
await prereg 080 (harness
\texttt{experiments/go14\_transfer.py}: per-cell subadditivity with
block-local denominators, the interleaved-prefix set identity, the
$\mathcal F_0$-only zero-claim with its $U$-coupled counter-values,
the (H$^\ast$) refutation netted as a gated fact, the optimizer
charges at $m\in\{8,16\}$, the constants and the corollary
arithmetic including the failing $n=24$ base, and reproduction of the
transfer anchors and of the within-class ladder). The scope proposed
for 080 is: Theorem~\ref{thm:transfer} and
Corollary~\ref{cor:lower} (unconditional), the transfer anchors, the
within-class ladder, and the conditional
Theorem~\ref{thm:transfer2} / Corollary~\ref{cor:plateau}.

The v0.6 material --- Theorem~\ref{thm:K} and its corollaries
(\S\ref{sec:K}), the H-repair lemma chain
L1--L4 and Theorem~\ref{thm:hrep}, the split-causal sub-family with
its certified $\phi^{\mathrm{sc}}$ brackets and
Theorem~\ref{thm:scstep}, the $\Pi$-retraction identities
(Lemma~\ref{lem:pi}), the extended monotone $m\times\Delta$ table, the
convexity-of-$\phi_m(D)$ hygiene of Remark~\ref{rem:feas}, the
recomputed corollary arithmetic under (H$^{\ast\ast}$) and under
(D)$+$(F), and the (H$^{\ast\ast\ast}$) value gaps --- awaits
\textbf{registration 081}, harness
\texttt{experiments/go14\_hrepair.py} (s1 the L1 identity and its
$m$-independence; s2 the L2 bound with its tightness witness; s3
L3/L4 and the H-repair inequality over a 22-cell pinned battery
including both refuting counterexamples; s4 the split-causal
sub-family, its constants and the linearity of the section; s5 the
$\Pi$-retraction identities; sK the three Theorem-K identities, the
two structural facts, the Cholesky split, $A_y$'s transpose symmetry
and the $m$-uniform bounds; s6 the block-distortion straddle and the
value-function convexity; s7 the constants and both corollary
arithmetics with the base-24 negative control; s8 the monotone
$m$-table; s9 $\phi^{\mathrm{sc}}_{2m}\le\phi_m$ and the
(H$^{\ast\ast\ast}$) gaps). Nothing in \S\ref{sec:K} or
\S\ref{sec:repair} is sealed.
\end{document}
